% ==============================================================================
% Common Preamble for IT Ethics Lectures
% ==============================================================================
% This file contains shared configuration for all lecture files.
% Include this after \documentclass using: % ==============================================================================
% Common Preamble for IT Ethics Lectures
% ==============================================================================
% This file contains shared configuration for all lecture files.
% Include this after \documentclass using: % ==============================================================================
% Common Preamble for IT Ethics Lectures
% ==============================================================================
% This file contains shared configuration for all lecture files.
% Include this after \documentclass using: % ==============================================================================
% Common Preamble for IT Ethics Lectures
% ==============================================================================
% This file contains shared configuration for all lecture files.
% Include this after \documentclass using: \input{lecture_preamble.tex}
%
% Individual lecture files should:
% 1. Declare \documentclass (beamer or beamerarticle)
% 2. \input{lecture_preamble.tex}
% 3. Define lecture-specific colors/customizations
% 4. Set title information
% ==============================================================================

% ------------------------------------------------------------------------------
% Theme and Basic Setup
% ------------------------------------------------------------------------------
\usetheme[progressbar=foot, block=fill, sectionpage=progressbar]{metropolis}

% ------------------------------------------------------------------------------
% Core Packages
% ------------------------------------------------------------------------------
\usepackage{tikz}
\usetikzlibrary{shapes,arrows,positioning,calc,decorations.pathreplacing,fit,backgrounds,chains}
\usepackage{booktabs}
\usepackage{array}
\usepackage{graphicx}
\usepackage{xcolor}
\usepackage[sfdefault]{plex-sans}   % IBM Plex Sans — compact, modern, tech-themed
\usepackage[T1]{fontenc}
\usepackage{fontawesome5}          % real vector icons

% Conditional text boxes (more flexible than basic beamer blocks)
\usepackage{tcolorbox}
\tcbuselibrary{skins,breakable}

% ------------------------------------------------------------------------------
% Bibliography (loaded in all modes for citations)
% ------------------------------------------------------------------------------
\usepackage[style=authoryear,backend=biber,natbib=true,maxcitenames=2]{biblatex}
\addbibresource{refs.bib}

% ------------------------------------------------------------------------------
% Article Mode Adjustments
% ------------------------------------------------------------------------------
% When compiling as article (using beamerarticle class), apply these settings
\mode<article>{
    % Use standard article sections instead of beamer frames
    \usepackage{fullpage}
    \usepackage{hyperref}
    \hypersetup{
        colorlinks=true,
        linkcolor=blue,
        citecolor=blue,
        urlcolor=blue
    }

    % Redefine frame environment for article mode
    % Frames become subsections in article mode
    \renewenvironment{frame}[1][]{
        \par\medskip
        \noindent\textbf{#1}\par\medskip
    }{}
}

% Presentation mode specific settings
\mode<presentation>{
    \setlength{\parskip}{0pt}

    % Show a TOC slide at the start of each section, current section highlighted
    \AtBeginSection[]{
        \begin{frame}{Lecture Outline}
            \tableofcontents[currentsection, hideallsubsections]
        \end{frame}
    }

    % FiraSans has ~15% taller metrics than Computer Modern; pull line spacing back
    % to avoid the massive overflows Metropolis + FiraSans produce by default.
    \renewcommand{\baselinestretch}{0.88}

    % Tighter list spacing
    \setbeamertemplate{itemize/enumerate body begin}{\setlength{\itemsep}{2pt}\setlength{\parsep}{0pt}}
    \setbeamertemplate{itemize/enumerate subbody begin}{\setlength{\itemsep}{1pt}\setlength{\parsep}{0pt}}
}

% ------------------------------------------------------------------------------
% Standard Colors for Boxes
% ------------------------------------------------------------------------------
% Consistent color scheme across all lectures for semantic elements
\definecolor{boxblue}{HTML}{1A4A7A}       % Arguments, concepts
\definecolor{boxgreen}{HTML}{2E7D32}      % Pro arguments, supporting points
\definecolor{boxred}{HTML}{C0392B}        % Objections, warnings
\definecolor{boxgold}{HTML}{B7770D}       % Case studies, examples
\definecolor{boxgray}{HTML}{555F6D}       % Quotes, citations
\definecolor{boxpurple}{HTML}{6A3093}     % Secondary accent
% Metropolis structural colors
\definecolor{mPrimary}{HTML}{1E3A5F}      % Deep navy — frametitle / title bg
\definecolor{mAccent}{HTML}{D4880A}       % Warm amber — progress bar / separators

% ------------------------------------------------------------------------------
% Beamer Theme Color Setup
% ------------------------------------------------------------------------------
% Applied once here; individual lectures should NOT override these.
\mode<presentation>{
    % Metropolis palette primary drives frametitle + title frame background
    \setbeamercolor{palette primary}{bg=mPrimary, fg=white}
    \setbeamercolor{frametitle}{bg=mPrimary, fg=white}
    % Title page element colors — must be explicit so they stay white
    % when a \usebackgroundtemplate image replaces the default navy fill
    \setbeamercolor{title}{fg=white}
    \setbeamercolor{subtitle}{fg=mAccent}
    \setbeamercolor{author}{fg=white}
    \setbeamercolor{institute}{fg=white!85!black}
    \setbeamercolor{date}{fg=white!75!black}
    % Progress bar and title separator use accent color
    \setbeamercolor{progress bar}{fg=mAccent, bg=mAccent!25}
    \setbeamercolor{title separator}{fg=mAccent}
    \setbeamercolor{alerted text}{fg=boxred}
    % Standard beamer blocks
    \setbeamercolor{block title}{fg=white, bg=boxblue}
    \setbeamercolor{block body}{bg=boxblue!8}
    \setbeamercolor{block title alerted}{fg=white, bg=boxred}
    \setbeamercolor{block body alerted}{bg=boxred!8}
    \setbeamercolor{block title example}{fg=white, bg=boxgreen}
    \setbeamercolor{block body example}{bg=boxgreen!5}
}

% ------------------------------------------------------------------------------
% Standard Custom Boxes
% ------------------------------------------------------------------------------
% Uniform boxes used across all lectures - do not redefine in individual files

% Argument/Reasoning Box (logical arguments in standard form)
\newtcolorbox{argumentbox}[1][Argument]{
    colback=boxblue!5,
    colframe=boxblue,
    fonttitle=\bfseries,
    title={#1},
    boxrule=1pt,
    arc=2pt
}

% Quotation Box (for philosophical quotes, sources)
\newtcolorbox{quotebox}[1][Quote]{
    colback=gray!10,
    colframe=boxgray,
    fonttitle=\itshape,
    title={#1},
    boxrule=0.5pt,
    arc=2pt,
    left=10pt,
    right=10pt,
    leftrule=3pt
}

% Case Study/Example Box
\newtcolorbox{casestudybox}[1][Case Study]{
    colback=boxgold!10,
    colframe=boxgold,
    fonttitle=\bfseries,
    title={#1},
    boxrule=1pt,
    arc=2pt
}

% Objection/Counterargument Box
\newtcolorbox{objectionbox}[1][Objection]{
    colback=boxred!5,
    colframe=boxred,
    fonttitle=\bfseries,
    title={#1},
    boxrule=1pt,
    arc=2pt
}

% Pro/Supporting Argument Box
\newtcolorbox{probox}[1][Supporting Argument]{
    colback=boxgreen!5,
    colframe=boxgreen,
    fonttitle=\bfseries,
    title={#1},
    boxrule=1pt,
    arc=2pt
}

% Concept/Definition Box
\newtcolorbox{conceptbox}[1][Key Concept]{
    colback=boxblue!5,
    colframe=boxblue,
    fonttitle=\bfseries,
    title={#1},
    boxrule=1pt,
    arc=2pt
}

% Discussion Box (for group activities, questions)
\newtcolorbox{discussionbox}[1][Discussion]{
    colback=boxgold!5,
    colframe=boxgold,
    fonttitle=\bfseries,
    title={#1},
    boxrule=1pt,
    arc=2pt
}

% Figure Box (for diagrams with explanations)
\newtcolorbox{figurebox}[1][Figure]{
    colback=boxblue!3,
    colframe=boxblue!70,
    fonttitle=\bfseries,
    title={#1},
    boxrule=0.75pt,
    arc=2pt
}

% ------------------------------------------------------------------------------
% Utility Commands and Icons
% ------------------------------------------------------------------------------

% fontawesome5 icon shorthands (\faXxx commands come from the fontawesome5 package)
\newcommand{\iconQuote}{\faQuoteLeft}
\newcommand{\iconBalance}{\faBalanceScale}
\newcommand{\iconDiscussion}{\faComments}
\newcommand{\iconIdea}{\faLightbulb}
\newcommand{\iconPro}{\faThumbsUp}
\newcommand{\iconCon}{\faThumbsDown}

% Discussion question command (for interactive moments)
\newcommand{\discussionquestion}[1]{%
    \begin{alertblock}{Discussion Question}
    #1
    \end{alertblock}
}

% fontawesome5 defines \faQuoteLeft, \faBalanceScale, etc. natively --- no aliases needed.
\newcommand{\discussion}[1]{\discussionquestion{#1}}

% Simple divider for slide sections
\newcommand{\sectiondivider}{%
    \begin{center}
    \rule{0.8\textwidth}{0.4pt}
    \end{center}
}

% Argument environment (for standard form arguments)
\newenvironment{argument}{%
  \begin{argumentbox}[Argument in Standard Form]
  \begin{enumerate}
}{%
  \end{enumerate}
  \end{argumentbox}
}

% Quote slide environment (combines icon and quotebox)
\newenvironment{quoteslide}[2]{%
  \begin{quotebox}[#1]
  \small
  \textit{``#2''}
}{%
  \end{quotebox}
}

% ==============================================================================
% End of Common Preamble
% ==============================================================================

%
% Individual lecture files should:
% 1. Declare \documentclass (beamer or beamerarticle)
% 2. % ==============================================================================
% Common Preamble for IT Ethics Lectures
% ==============================================================================
% This file contains shared configuration for all lecture files.
% Include this after \documentclass using: \input{lecture_preamble.tex}
%
% Individual lecture files should:
% 1. Declare \documentclass (beamer or beamerarticle)
% 2. \input{lecture_preamble.tex}
% 3. Define lecture-specific colors/customizations
% 4. Set title information
% ==============================================================================

% ------------------------------------------------------------------------------
% Theme and Basic Setup
% ------------------------------------------------------------------------------
\usetheme[progressbar=foot, block=fill, sectionpage=progressbar]{metropolis}

% ------------------------------------------------------------------------------
% Core Packages
% ------------------------------------------------------------------------------
\usepackage{tikz}
\usetikzlibrary{shapes,arrows,positioning,calc,decorations.pathreplacing,fit,backgrounds,chains}
\usepackage{booktabs}
\usepackage{array}
\usepackage{graphicx}
\usepackage{xcolor}
\usepackage[sfdefault]{plex-sans}   % IBM Plex Sans — compact, modern, tech-themed
\usepackage[T1]{fontenc}
\usepackage{fontawesome5}          % real vector icons

% Conditional text boxes (more flexible than basic beamer blocks)
\usepackage{tcolorbox}
\tcbuselibrary{skins,breakable}

% ------------------------------------------------------------------------------
% Bibliography (loaded in all modes for citations)
% ------------------------------------------------------------------------------
\usepackage[style=authoryear,backend=biber,natbib=true,maxcitenames=2]{biblatex}
\addbibresource{refs.bib}

% ------------------------------------------------------------------------------
% Article Mode Adjustments
% ------------------------------------------------------------------------------
% When compiling as article (using beamerarticle class), apply these settings
\mode<article>{
    % Use standard article sections instead of beamer frames
    \usepackage{fullpage}
    \usepackage{hyperref}
    \hypersetup{
        colorlinks=true,
        linkcolor=blue,
        citecolor=blue,
        urlcolor=blue
    }

    % Redefine frame environment for article mode
    % Frames become subsections in article mode
    \renewenvironment{frame}[1][]{
        \par\medskip
        \noindent\textbf{#1}\par\medskip
    }{}
}

% Presentation mode specific settings
\mode<presentation>{
    \setlength{\parskip}{0pt}

    % Show a TOC slide at the start of each section, current section highlighted
    \AtBeginSection[]{
        \begin{frame}{Lecture Outline}
            \tableofcontents[currentsection, hideallsubsections]
        \end{frame}
    }

    % FiraSans has ~15% taller metrics than Computer Modern; pull line spacing back
    % to avoid the massive overflows Metropolis + FiraSans produce by default.
    \renewcommand{\baselinestretch}{0.88}

    % Tighter list spacing
    \setbeamertemplate{itemize/enumerate body begin}{\setlength{\itemsep}{2pt}\setlength{\parsep}{0pt}}
    \setbeamertemplate{itemize/enumerate subbody begin}{\setlength{\itemsep}{1pt}\setlength{\parsep}{0pt}}
}

% ------------------------------------------------------------------------------
% Standard Colors for Boxes
% ------------------------------------------------------------------------------
% Consistent color scheme across all lectures for semantic elements
\definecolor{boxblue}{HTML}{1A4A7A}       % Arguments, concepts
\definecolor{boxgreen}{HTML}{2E7D32}      % Pro arguments, supporting points
\definecolor{boxred}{HTML}{C0392B}        % Objections, warnings
\definecolor{boxgold}{HTML}{B7770D}       % Case studies, examples
\definecolor{boxgray}{HTML}{555F6D}       % Quotes, citations
\definecolor{boxpurple}{HTML}{6A3093}     % Secondary accent
% Metropolis structural colors
\definecolor{mPrimary}{HTML}{1E3A5F}      % Deep navy — frametitle / title bg
\definecolor{mAccent}{HTML}{D4880A}       % Warm amber — progress bar / separators

% ------------------------------------------------------------------------------
% Beamer Theme Color Setup
% ------------------------------------------------------------------------------
% Applied once here; individual lectures should NOT override these.
\mode<presentation>{
    % Metropolis palette primary drives frametitle + title frame background
    \setbeamercolor{palette primary}{bg=mPrimary, fg=white}
    \setbeamercolor{frametitle}{bg=mPrimary, fg=white}
    % Title page element colors — must be explicit so they stay white
    % when a \usebackgroundtemplate image replaces the default navy fill
    \setbeamercolor{title}{fg=white}
    \setbeamercolor{subtitle}{fg=mAccent}
    \setbeamercolor{author}{fg=white}
    \setbeamercolor{institute}{fg=white!85!black}
    \setbeamercolor{date}{fg=white!75!black}
    % Progress bar and title separator use accent color
    \setbeamercolor{progress bar}{fg=mAccent, bg=mAccent!25}
    \setbeamercolor{title separator}{fg=mAccent}
    \setbeamercolor{alerted text}{fg=boxred}
    % Standard beamer blocks
    \setbeamercolor{block title}{fg=white, bg=boxblue}
    \setbeamercolor{block body}{bg=boxblue!8}
    \setbeamercolor{block title alerted}{fg=white, bg=boxred}
    \setbeamercolor{block body alerted}{bg=boxred!8}
    \setbeamercolor{block title example}{fg=white, bg=boxgreen}
    \setbeamercolor{block body example}{bg=boxgreen!5}
}

% ------------------------------------------------------------------------------
% Standard Custom Boxes
% ------------------------------------------------------------------------------
% Uniform boxes used across all lectures - do not redefine in individual files

% Argument/Reasoning Box (logical arguments in standard form)
\newtcolorbox{argumentbox}[1][Argument]{
    colback=boxblue!5,
    colframe=boxblue,
    fonttitle=\bfseries,
    title={#1},
    boxrule=1pt,
    arc=2pt
}

% Quotation Box (for philosophical quotes, sources)
\newtcolorbox{quotebox}[1][Quote]{
    colback=gray!10,
    colframe=boxgray,
    fonttitle=\itshape,
    title={#1},
    boxrule=0.5pt,
    arc=2pt,
    left=10pt,
    right=10pt,
    leftrule=3pt
}

% Case Study/Example Box
\newtcolorbox{casestudybox}[1][Case Study]{
    colback=boxgold!10,
    colframe=boxgold,
    fonttitle=\bfseries,
    title={#1},
    boxrule=1pt,
    arc=2pt
}

% Objection/Counterargument Box
\newtcolorbox{objectionbox}[1][Objection]{
    colback=boxred!5,
    colframe=boxred,
    fonttitle=\bfseries,
    title={#1},
    boxrule=1pt,
    arc=2pt
}

% Pro/Supporting Argument Box
\newtcolorbox{probox}[1][Supporting Argument]{
    colback=boxgreen!5,
    colframe=boxgreen,
    fonttitle=\bfseries,
    title={#1},
    boxrule=1pt,
    arc=2pt
}

% Concept/Definition Box
\newtcolorbox{conceptbox}[1][Key Concept]{
    colback=boxblue!5,
    colframe=boxblue,
    fonttitle=\bfseries,
    title={#1},
    boxrule=1pt,
    arc=2pt
}

% Discussion Box (for group activities, questions)
\newtcolorbox{discussionbox}[1][Discussion]{
    colback=boxgold!5,
    colframe=boxgold,
    fonttitle=\bfseries,
    title={#1},
    boxrule=1pt,
    arc=2pt
}

% Figure Box (for diagrams with explanations)
\newtcolorbox{figurebox}[1][Figure]{
    colback=boxblue!3,
    colframe=boxblue!70,
    fonttitle=\bfseries,
    title={#1},
    boxrule=0.75pt,
    arc=2pt
}

% ------------------------------------------------------------------------------
% Utility Commands and Icons
% ------------------------------------------------------------------------------

% fontawesome5 icon shorthands (\faXxx commands come from the fontawesome5 package)
\newcommand{\iconQuote}{\faQuoteLeft}
\newcommand{\iconBalance}{\faBalanceScale}
\newcommand{\iconDiscussion}{\faComments}
\newcommand{\iconIdea}{\faLightbulb}
\newcommand{\iconPro}{\faThumbsUp}
\newcommand{\iconCon}{\faThumbsDown}

% Discussion question command (for interactive moments)
\newcommand{\discussionquestion}[1]{%
    \begin{alertblock}{Discussion Question}
    #1
    \end{alertblock}
}

% fontawesome5 defines \faQuoteLeft, \faBalanceScale, etc. natively --- no aliases needed.
\newcommand{\discussion}[1]{\discussionquestion{#1}}

% Simple divider for slide sections
\newcommand{\sectiondivider}{%
    \begin{center}
    \rule{0.8\textwidth}{0.4pt}
    \end{center}
}

% Argument environment (for standard form arguments)
\newenvironment{argument}{%
  \begin{argumentbox}[Argument in Standard Form]
  \begin{enumerate}
}{%
  \end{enumerate}
  \end{argumentbox}
}

% Quote slide environment (combines icon and quotebox)
\newenvironment{quoteslide}[2]{%
  \begin{quotebox}[#1]
  \small
  \textit{``#2''}
}{%
  \end{quotebox}
}

% ==============================================================================
% End of Common Preamble
% ==============================================================================

% 3. Define lecture-specific colors/customizations
% 4. Set title information
% ==============================================================================

% ------------------------------------------------------------------------------
% Theme and Basic Setup
% ------------------------------------------------------------------------------
\usetheme[progressbar=foot, block=fill, sectionpage=progressbar]{metropolis}

% ------------------------------------------------------------------------------
% Core Packages
% ------------------------------------------------------------------------------
\usepackage{tikz}
\usetikzlibrary{shapes,arrows,positioning,calc,decorations.pathreplacing,fit,backgrounds,chains}
\usepackage{booktabs}
\usepackage{array}
\usepackage{graphicx}
\usepackage{xcolor}
\usepackage[sfdefault]{plex-sans}   % IBM Plex Sans — compact, modern, tech-themed
\usepackage[T1]{fontenc}
\usepackage{fontawesome5}          % real vector icons

% Conditional text boxes (more flexible than basic beamer blocks)
\usepackage{tcolorbox}
\tcbuselibrary{skins,breakable}

% ------------------------------------------------------------------------------
% Bibliography (loaded in all modes for citations)
% ------------------------------------------------------------------------------
\usepackage[style=authoryear,backend=biber,natbib=true,maxcitenames=2]{biblatex}
\addbibresource{refs.bib}

% ------------------------------------------------------------------------------
% Article Mode Adjustments
% ------------------------------------------------------------------------------
% When compiling as article (using beamerarticle class), apply these settings
\mode<article>{
    % Use standard article sections instead of beamer frames
    \usepackage{fullpage}
    \usepackage{hyperref}
    \hypersetup{
        colorlinks=true,
        linkcolor=blue,
        citecolor=blue,
        urlcolor=blue
    }

    % Redefine frame environment for article mode
    % Frames become subsections in article mode
    \renewenvironment{frame}[1][]{
        \par\medskip
        \noindent\textbf{#1}\par\medskip
    }{}
}

% Presentation mode specific settings
\mode<presentation>{
    \setlength{\parskip}{0pt}

    % Show a TOC slide at the start of each section, current section highlighted
    \AtBeginSection[]{
        \begin{frame}{Lecture Outline}
            \tableofcontents[currentsection, hideallsubsections]
        \end{frame}
    }

    % FiraSans has ~15% taller metrics than Computer Modern; pull line spacing back
    % to avoid the massive overflows Metropolis + FiraSans produce by default.
    \renewcommand{\baselinestretch}{0.88}

    % Tighter list spacing
    \setbeamertemplate{itemize/enumerate body begin}{\setlength{\itemsep}{2pt}\setlength{\parsep}{0pt}}
    \setbeamertemplate{itemize/enumerate subbody begin}{\setlength{\itemsep}{1pt}\setlength{\parsep}{0pt}}
}

% ------------------------------------------------------------------------------
% Standard Colors for Boxes
% ------------------------------------------------------------------------------
% Consistent color scheme across all lectures for semantic elements
\definecolor{boxblue}{HTML}{1A4A7A}       % Arguments, concepts
\definecolor{boxgreen}{HTML}{2E7D32}      % Pro arguments, supporting points
\definecolor{boxred}{HTML}{C0392B}        % Objections, warnings
\definecolor{boxgold}{HTML}{B7770D}       % Case studies, examples
\definecolor{boxgray}{HTML}{555F6D}       % Quotes, citations
\definecolor{boxpurple}{HTML}{6A3093}     % Secondary accent
% Metropolis structural colors
\definecolor{mPrimary}{HTML}{1E3A5F}      % Deep navy — frametitle / title bg
\definecolor{mAccent}{HTML}{D4880A}       % Warm amber — progress bar / separators

% ------------------------------------------------------------------------------
% Beamer Theme Color Setup
% ------------------------------------------------------------------------------
% Applied once here; individual lectures should NOT override these.
\mode<presentation>{
    % Metropolis palette primary drives frametitle + title frame background
    \setbeamercolor{palette primary}{bg=mPrimary, fg=white}
    \setbeamercolor{frametitle}{bg=mPrimary, fg=white}
    % Title page element colors — must be explicit so they stay white
    % when a \usebackgroundtemplate image replaces the default navy fill
    \setbeamercolor{title}{fg=white}
    \setbeamercolor{subtitle}{fg=mAccent}
    \setbeamercolor{author}{fg=white}
    \setbeamercolor{institute}{fg=white!85!black}
    \setbeamercolor{date}{fg=white!75!black}
    % Progress bar and title separator use accent color
    \setbeamercolor{progress bar}{fg=mAccent, bg=mAccent!25}
    \setbeamercolor{title separator}{fg=mAccent}
    \setbeamercolor{alerted text}{fg=boxred}
    % Standard beamer blocks
    \setbeamercolor{block title}{fg=white, bg=boxblue}
    \setbeamercolor{block body}{bg=boxblue!8}
    \setbeamercolor{block title alerted}{fg=white, bg=boxred}
    \setbeamercolor{block body alerted}{bg=boxred!8}
    \setbeamercolor{block title example}{fg=white, bg=boxgreen}
    \setbeamercolor{block body example}{bg=boxgreen!5}
}

% ------------------------------------------------------------------------------
% Standard Custom Boxes
% ------------------------------------------------------------------------------
% Uniform boxes used across all lectures - do not redefine in individual files

% Argument/Reasoning Box (logical arguments in standard form)
\newtcolorbox{argumentbox}[1][Argument]{
    colback=boxblue!5,
    colframe=boxblue,
    fonttitle=\bfseries,
    title={#1},
    boxrule=1pt,
    arc=2pt
}

% Quotation Box (for philosophical quotes, sources)
\newtcolorbox{quotebox}[1][Quote]{
    colback=gray!10,
    colframe=boxgray,
    fonttitle=\itshape,
    title={#1},
    boxrule=0.5pt,
    arc=2pt,
    left=10pt,
    right=10pt,
    leftrule=3pt
}

% Case Study/Example Box
\newtcolorbox{casestudybox}[1][Case Study]{
    colback=boxgold!10,
    colframe=boxgold,
    fonttitle=\bfseries,
    title={#1},
    boxrule=1pt,
    arc=2pt
}

% Objection/Counterargument Box
\newtcolorbox{objectionbox}[1][Objection]{
    colback=boxred!5,
    colframe=boxred,
    fonttitle=\bfseries,
    title={#1},
    boxrule=1pt,
    arc=2pt
}

% Pro/Supporting Argument Box
\newtcolorbox{probox}[1][Supporting Argument]{
    colback=boxgreen!5,
    colframe=boxgreen,
    fonttitle=\bfseries,
    title={#1},
    boxrule=1pt,
    arc=2pt
}

% Concept/Definition Box
\newtcolorbox{conceptbox}[1][Key Concept]{
    colback=boxblue!5,
    colframe=boxblue,
    fonttitle=\bfseries,
    title={#1},
    boxrule=1pt,
    arc=2pt
}

% Discussion Box (for group activities, questions)
\newtcolorbox{discussionbox}[1][Discussion]{
    colback=boxgold!5,
    colframe=boxgold,
    fonttitle=\bfseries,
    title={#1},
    boxrule=1pt,
    arc=2pt
}

% Figure Box (for diagrams with explanations)
\newtcolorbox{figurebox}[1][Figure]{
    colback=boxblue!3,
    colframe=boxblue!70,
    fonttitle=\bfseries,
    title={#1},
    boxrule=0.75pt,
    arc=2pt
}

% ------------------------------------------------------------------------------
% Utility Commands and Icons
% ------------------------------------------------------------------------------

% fontawesome5 icon shorthands (\faXxx commands come from the fontawesome5 package)
\newcommand{\iconQuote}{\faQuoteLeft}
\newcommand{\iconBalance}{\faBalanceScale}
\newcommand{\iconDiscussion}{\faComments}
\newcommand{\iconIdea}{\faLightbulb}
\newcommand{\iconPro}{\faThumbsUp}
\newcommand{\iconCon}{\faThumbsDown}

% Discussion question command (for interactive moments)
\newcommand{\discussionquestion}[1]{%
    \begin{alertblock}{Discussion Question}
    #1
    \end{alertblock}
}

% fontawesome5 defines \faQuoteLeft, \faBalanceScale, etc. natively --- no aliases needed.
\newcommand{\discussion}[1]{\discussionquestion{#1}}

% Simple divider for slide sections
\newcommand{\sectiondivider}{%
    \begin{center}
    \rule{0.8\textwidth}{0.4pt}
    \end{center}
}

% Argument environment (for standard form arguments)
\newenvironment{argument}{%
  \begin{argumentbox}[Argument in Standard Form]
  \begin{enumerate}
}{%
  \end{enumerate}
  \end{argumentbox}
}

% Quote slide environment (combines icon and quotebox)
\newenvironment{quoteslide}[2]{%
  \begin{quotebox}[#1]
  \small
  \textit{``#2''}
}{%
  \end{quotebox}
}

% ==============================================================================
% End of Common Preamble
% ==============================================================================

%
% Individual lecture files should:
% 1. Declare \documentclass (beamer or beamerarticle)
% 2. % ==============================================================================
% Common Preamble for IT Ethics Lectures
% ==============================================================================
% This file contains shared configuration for all lecture files.
% Include this after \documentclass using: % ==============================================================================
% Common Preamble for IT Ethics Lectures
% ==============================================================================
% This file contains shared configuration for all lecture files.
% Include this after \documentclass using: \input{lecture_preamble.tex}
%
% Individual lecture files should:
% 1. Declare \documentclass (beamer or beamerarticle)
% 2. \input{lecture_preamble.tex}
% 3. Define lecture-specific colors/customizations
% 4. Set title information
% ==============================================================================

% ------------------------------------------------------------------------------
% Theme and Basic Setup
% ------------------------------------------------------------------------------
\usetheme[progressbar=foot, block=fill, sectionpage=progressbar]{metropolis}

% ------------------------------------------------------------------------------
% Core Packages
% ------------------------------------------------------------------------------
\usepackage{tikz}
\usetikzlibrary{shapes,arrows,positioning,calc,decorations.pathreplacing,fit,backgrounds,chains}
\usepackage{booktabs}
\usepackage{array}
\usepackage{graphicx}
\usepackage{xcolor}
\usepackage[sfdefault]{plex-sans}   % IBM Plex Sans — compact, modern, tech-themed
\usepackage[T1]{fontenc}
\usepackage{fontawesome5}          % real vector icons

% Conditional text boxes (more flexible than basic beamer blocks)
\usepackage{tcolorbox}
\tcbuselibrary{skins,breakable}

% ------------------------------------------------------------------------------
% Bibliography (loaded in all modes for citations)
% ------------------------------------------------------------------------------
\usepackage[style=authoryear,backend=biber,natbib=true,maxcitenames=2]{biblatex}
\addbibresource{refs.bib}

% ------------------------------------------------------------------------------
% Article Mode Adjustments
% ------------------------------------------------------------------------------
% When compiling as article (using beamerarticle class), apply these settings
\mode<article>{
    % Use standard article sections instead of beamer frames
    \usepackage{fullpage}
    \usepackage{hyperref}
    \hypersetup{
        colorlinks=true,
        linkcolor=blue,
        citecolor=blue,
        urlcolor=blue
    }

    % Redefine frame environment for article mode
    % Frames become subsections in article mode
    \renewenvironment{frame}[1][]{
        \par\medskip
        \noindent\textbf{#1}\par\medskip
    }{}
}

% Presentation mode specific settings
\mode<presentation>{
    \setlength{\parskip}{0pt}

    % Show a TOC slide at the start of each section, current section highlighted
    \AtBeginSection[]{
        \begin{frame}{Lecture Outline}
            \tableofcontents[currentsection, hideallsubsections]
        \end{frame}
    }

    % FiraSans has ~15% taller metrics than Computer Modern; pull line spacing back
    % to avoid the massive overflows Metropolis + FiraSans produce by default.
    \renewcommand{\baselinestretch}{0.88}

    % Tighter list spacing
    \setbeamertemplate{itemize/enumerate body begin}{\setlength{\itemsep}{2pt}\setlength{\parsep}{0pt}}
    \setbeamertemplate{itemize/enumerate subbody begin}{\setlength{\itemsep}{1pt}\setlength{\parsep}{0pt}}
}

% ------------------------------------------------------------------------------
% Standard Colors for Boxes
% ------------------------------------------------------------------------------
% Consistent color scheme across all lectures for semantic elements
\definecolor{boxblue}{HTML}{1A4A7A}       % Arguments, concepts
\definecolor{boxgreen}{HTML}{2E7D32}      % Pro arguments, supporting points
\definecolor{boxred}{HTML}{C0392B}        % Objections, warnings
\definecolor{boxgold}{HTML}{B7770D}       % Case studies, examples
\definecolor{boxgray}{HTML}{555F6D}       % Quotes, citations
\definecolor{boxpurple}{HTML}{6A3093}     % Secondary accent
% Metropolis structural colors
\definecolor{mPrimary}{HTML}{1E3A5F}      % Deep navy — frametitle / title bg
\definecolor{mAccent}{HTML}{D4880A}       % Warm amber — progress bar / separators

% ------------------------------------------------------------------------------
% Beamer Theme Color Setup
% ------------------------------------------------------------------------------
% Applied once here; individual lectures should NOT override these.
\mode<presentation>{
    % Metropolis palette primary drives frametitle + title frame background
    \setbeamercolor{palette primary}{bg=mPrimary, fg=white}
    \setbeamercolor{frametitle}{bg=mPrimary, fg=white}
    % Title page element colors — must be explicit so they stay white
    % when a \usebackgroundtemplate image replaces the default navy fill
    \setbeamercolor{title}{fg=white}
    \setbeamercolor{subtitle}{fg=mAccent}
    \setbeamercolor{author}{fg=white}
    \setbeamercolor{institute}{fg=white!85!black}
    \setbeamercolor{date}{fg=white!75!black}
    % Progress bar and title separator use accent color
    \setbeamercolor{progress bar}{fg=mAccent, bg=mAccent!25}
    \setbeamercolor{title separator}{fg=mAccent}
    \setbeamercolor{alerted text}{fg=boxred}
    % Standard beamer blocks
    \setbeamercolor{block title}{fg=white, bg=boxblue}
    \setbeamercolor{block body}{bg=boxblue!8}
    \setbeamercolor{block title alerted}{fg=white, bg=boxred}
    \setbeamercolor{block body alerted}{bg=boxred!8}
    \setbeamercolor{block title example}{fg=white, bg=boxgreen}
    \setbeamercolor{block body example}{bg=boxgreen!5}
}

% ------------------------------------------------------------------------------
% Standard Custom Boxes
% ------------------------------------------------------------------------------
% Uniform boxes used across all lectures - do not redefine in individual files

% Argument/Reasoning Box (logical arguments in standard form)
\newtcolorbox{argumentbox}[1][Argument]{
    colback=boxblue!5,
    colframe=boxblue,
    fonttitle=\bfseries,
    title={#1},
    boxrule=1pt,
    arc=2pt
}

% Quotation Box (for philosophical quotes, sources)
\newtcolorbox{quotebox}[1][Quote]{
    colback=gray!10,
    colframe=boxgray,
    fonttitle=\itshape,
    title={#1},
    boxrule=0.5pt,
    arc=2pt,
    left=10pt,
    right=10pt,
    leftrule=3pt
}

% Case Study/Example Box
\newtcolorbox{casestudybox}[1][Case Study]{
    colback=boxgold!10,
    colframe=boxgold,
    fonttitle=\bfseries,
    title={#1},
    boxrule=1pt,
    arc=2pt
}

% Objection/Counterargument Box
\newtcolorbox{objectionbox}[1][Objection]{
    colback=boxred!5,
    colframe=boxred,
    fonttitle=\bfseries,
    title={#1},
    boxrule=1pt,
    arc=2pt
}

% Pro/Supporting Argument Box
\newtcolorbox{probox}[1][Supporting Argument]{
    colback=boxgreen!5,
    colframe=boxgreen,
    fonttitle=\bfseries,
    title={#1},
    boxrule=1pt,
    arc=2pt
}

% Concept/Definition Box
\newtcolorbox{conceptbox}[1][Key Concept]{
    colback=boxblue!5,
    colframe=boxblue,
    fonttitle=\bfseries,
    title={#1},
    boxrule=1pt,
    arc=2pt
}

% Discussion Box (for group activities, questions)
\newtcolorbox{discussionbox}[1][Discussion]{
    colback=boxgold!5,
    colframe=boxgold,
    fonttitle=\bfseries,
    title={#1},
    boxrule=1pt,
    arc=2pt
}

% Figure Box (for diagrams with explanations)
\newtcolorbox{figurebox}[1][Figure]{
    colback=boxblue!3,
    colframe=boxblue!70,
    fonttitle=\bfseries,
    title={#1},
    boxrule=0.75pt,
    arc=2pt
}

% ------------------------------------------------------------------------------
% Utility Commands and Icons
% ------------------------------------------------------------------------------

% fontawesome5 icon shorthands (\faXxx commands come from the fontawesome5 package)
\newcommand{\iconQuote}{\faQuoteLeft}
\newcommand{\iconBalance}{\faBalanceScale}
\newcommand{\iconDiscussion}{\faComments}
\newcommand{\iconIdea}{\faLightbulb}
\newcommand{\iconPro}{\faThumbsUp}
\newcommand{\iconCon}{\faThumbsDown}

% Discussion question command (for interactive moments)
\newcommand{\discussionquestion}[1]{%
    \begin{alertblock}{Discussion Question}
    #1
    \end{alertblock}
}

% fontawesome5 defines \faQuoteLeft, \faBalanceScale, etc. natively --- no aliases needed.
\newcommand{\discussion}[1]{\discussionquestion{#1}}

% Simple divider for slide sections
\newcommand{\sectiondivider}{%
    \begin{center}
    \rule{0.8\textwidth}{0.4pt}
    \end{center}
}

% Argument environment (for standard form arguments)
\newenvironment{argument}{%
  \begin{argumentbox}[Argument in Standard Form]
  \begin{enumerate}
}{%
  \end{enumerate}
  \end{argumentbox}
}

% Quote slide environment (combines icon and quotebox)
\newenvironment{quoteslide}[2]{%
  \begin{quotebox}[#1]
  \small
  \textit{``#2''}
}{%
  \end{quotebox}
}

% ==============================================================================
% End of Common Preamble
% ==============================================================================

%
% Individual lecture files should:
% 1. Declare \documentclass (beamer or beamerarticle)
% 2. % ==============================================================================
% Common Preamble for IT Ethics Lectures
% ==============================================================================
% This file contains shared configuration for all lecture files.
% Include this after \documentclass using: \input{lecture_preamble.tex}
%
% Individual lecture files should:
% 1. Declare \documentclass (beamer or beamerarticle)
% 2. \input{lecture_preamble.tex}
% 3. Define lecture-specific colors/customizations
% 4. Set title information
% ==============================================================================

% ------------------------------------------------------------------------------
% Theme and Basic Setup
% ------------------------------------------------------------------------------
\usetheme[progressbar=foot, block=fill, sectionpage=progressbar]{metropolis}

% ------------------------------------------------------------------------------
% Core Packages
% ------------------------------------------------------------------------------
\usepackage{tikz}
\usetikzlibrary{shapes,arrows,positioning,calc,decorations.pathreplacing,fit,backgrounds,chains}
\usepackage{booktabs}
\usepackage{array}
\usepackage{graphicx}
\usepackage{xcolor}
\usepackage[sfdefault]{plex-sans}   % IBM Plex Sans — compact, modern, tech-themed
\usepackage[T1]{fontenc}
\usepackage{fontawesome5}          % real vector icons

% Conditional text boxes (more flexible than basic beamer blocks)
\usepackage{tcolorbox}
\tcbuselibrary{skins,breakable}

% ------------------------------------------------------------------------------
% Bibliography (loaded in all modes for citations)
% ------------------------------------------------------------------------------
\usepackage[style=authoryear,backend=biber,natbib=true,maxcitenames=2]{biblatex}
\addbibresource{refs.bib}

% ------------------------------------------------------------------------------
% Article Mode Adjustments
% ------------------------------------------------------------------------------
% When compiling as article (using beamerarticle class), apply these settings
\mode<article>{
    % Use standard article sections instead of beamer frames
    \usepackage{fullpage}
    \usepackage{hyperref}
    \hypersetup{
        colorlinks=true,
        linkcolor=blue,
        citecolor=blue,
        urlcolor=blue
    }

    % Redefine frame environment for article mode
    % Frames become subsections in article mode
    \renewenvironment{frame}[1][]{
        \par\medskip
        \noindent\textbf{#1}\par\medskip
    }{}
}

% Presentation mode specific settings
\mode<presentation>{
    \setlength{\parskip}{0pt}

    % Show a TOC slide at the start of each section, current section highlighted
    \AtBeginSection[]{
        \begin{frame}{Lecture Outline}
            \tableofcontents[currentsection, hideallsubsections]
        \end{frame}
    }

    % FiraSans has ~15% taller metrics than Computer Modern; pull line spacing back
    % to avoid the massive overflows Metropolis + FiraSans produce by default.
    \renewcommand{\baselinestretch}{0.88}

    % Tighter list spacing
    \setbeamertemplate{itemize/enumerate body begin}{\setlength{\itemsep}{2pt}\setlength{\parsep}{0pt}}
    \setbeamertemplate{itemize/enumerate subbody begin}{\setlength{\itemsep}{1pt}\setlength{\parsep}{0pt}}
}

% ------------------------------------------------------------------------------
% Standard Colors for Boxes
% ------------------------------------------------------------------------------
% Consistent color scheme across all lectures for semantic elements
\definecolor{boxblue}{HTML}{1A4A7A}       % Arguments, concepts
\definecolor{boxgreen}{HTML}{2E7D32}      % Pro arguments, supporting points
\definecolor{boxred}{HTML}{C0392B}        % Objections, warnings
\definecolor{boxgold}{HTML}{B7770D}       % Case studies, examples
\definecolor{boxgray}{HTML}{555F6D}       % Quotes, citations
\definecolor{boxpurple}{HTML}{6A3093}     % Secondary accent
% Metropolis structural colors
\definecolor{mPrimary}{HTML}{1E3A5F}      % Deep navy — frametitle / title bg
\definecolor{mAccent}{HTML}{D4880A}       % Warm amber — progress bar / separators

% ------------------------------------------------------------------------------
% Beamer Theme Color Setup
% ------------------------------------------------------------------------------
% Applied once here; individual lectures should NOT override these.
\mode<presentation>{
    % Metropolis palette primary drives frametitle + title frame background
    \setbeamercolor{palette primary}{bg=mPrimary, fg=white}
    \setbeamercolor{frametitle}{bg=mPrimary, fg=white}
    % Title page element colors — must be explicit so they stay white
    % when a \usebackgroundtemplate image replaces the default navy fill
    \setbeamercolor{title}{fg=white}
    \setbeamercolor{subtitle}{fg=mAccent}
    \setbeamercolor{author}{fg=white}
    \setbeamercolor{institute}{fg=white!85!black}
    \setbeamercolor{date}{fg=white!75!black}
    % Progress bar and title separator use accent color
    \setbeamercolor{progress bar}{fg=mAccent, bg=mAccent!25}
    \setbeamercolor{title separator}{fg=mAccent}
    \setbeamercolor{alerted text}{fg=boxred}
    % Standard beamer blocks
    \setbeamercolor{block title}{fg=white, bg=boxblue}
    \setbeamercolor{block body}{bg=boxblue!8}
    \setbeamercolor{block title alerted}{fg=white, bg=boxred}
    \setbeamercolor{block body alerted}{bg=boxred!8}
    \setbeamercolor{block title example}{fg=white, bg=boxgreen}
    \setbeamercolor{block body example}{bg=boxgreen!5}
}

% ------------------------------------------------------------------------------
% Standard Custom Boxes
% ------------------------------------------------------------------------------
% Uniform boxes used across all lectures - do not redefine in individual files

% Argument/Reasoning Box (logical arguments in standard form)
\newtcolorbox{argumentbox}[1][Argument]{
    colback=boxblue!5,
    colframe=boxblue,
    fonttitle=\bfseries,
    title={#1},
    boxrule=1pt,
    arc=2pt
}

% Quotation Box (for philosophical quotes, sources)
\newtcolorbox{quotebox}[1][Quote]{
    colback=gray!10,
    colframe=boxgray,
    fonttitle=\itshape,
    title={#1},
    boxrule=0.5pt,
    arc=2pt,
    left=10pt,
    right=10pt,
    leftrule=3pt
}

% Case Study/Example Box
\newtcolorbox{casestudybox}[1][Case Study]{
    colback=boxgold!10,
    colframe=boxgold,
    fonttitle=\bfseries,
    title={#1},
    boxrule=1pt,
    arc=2pt
}

% Objection/Counterargument Box
\newtcolorbox{objectionbox}[1][Objection]{
    colback=boxred!5,
    colframe=boxred,
    fonttitle=\bfseries,
    title={#1},
    boxrule=1pt,
    arc=2pt
}

% Pro/Supporting Argument Box
\newtcolorbox{probox}[1][Supporting Argument]{
    colback=boxgreen!5,
    colframe=boxgreen,
    fonttitle=\bfseries,
    title={#1},
    boxrule=1pt,
    arc=2pt
}

% Concept/Definition Box
\newtcolorbox{conceptbox}[1][Key Concept]{
    colback=boxblue!5,
    colframe=boxblue,
    fonttitle=\bfseries,
    title={#1},
    boxrule=1pt,
    arc=2pt
}

% Discussion Box (for group activities, questions)
\newtcolorbox{discussionbox}[1][Discussion]{
    colback=boxgold!5,
    colframe=boxgold,
    fonttitle=\bfseries,
    title={#1},
    boxrule=1pt,
    arc=2pt
}

% Figure Box (for diagrams with explanations)
\newtcolorbox{figurebox}[1][Figure]{
    colback=boxblue!3,
    colframe=boxblue!70,
    fonttitle=\bfseries,
    title={#1},
    boxrule=0.75pt,
    arc=2pt
}

% ------------------------------------------------------------------------------
% Utility Commands and Icons
% ------------------------------------------------------------------------------

% fontawesome5 icon shorthands (\faXxx commands come from the fontawesome5 package)
\newcommand{\iconQuote}{\faQuoteLeft}
\newcommand{\iconBalance}{\faBalanceScale}
\newcommand{\iconDiscussion}{\faComments}
\newcommand{\iconIdea}{\faLightbulb}
\newcommand{\iconPro}{\faThumbsUp}
\newcommand{\iconCon}{\faThumbsDown}

% Discussion question command (for interactive moments)
\newcommand{\discussionquestion}[1]{%
    \begin{alertblock}{Discussion Question}
    #1
    \end{alertblock}
}

% fontawesome5 defines \faQuoteLeft, \faBalanceScale, etc. natively --- no aliases needed.
\newcommand{\discussion}[1]{\discussionquestion{#1}}

% Simple divider for slide sections
\newcommand{\sectiondivider}{%
    \begin{center}
    \rule{0.8\textwidth}{0.4pt}
    \end{center}
}

% Argument environment (for standard form arguments)
\newenvironment{argument}{%
  \begin{argumentbox}[Argument in Standard Form]
  \begin{enumerate}
}{%
  \end{enumerate}
  \end{argumentbox}
}

% Quote slide environment (combines icon and quotebox)
\newenvironment{quoteslide}[2]{%
  \begin{quotebox}[#1]
  \small
  \textit{``#2''}
}{%
  \end{quotebox}
}

% ==============================================================================
% End of Common Preamble
% ==============================================================================

% 3. Define lecture-specific colors/customizations
% 4. Set title information
% ==============================================================================

% ------------------------------------------------------------------------------
% Theme and Basic Setup
% ------------------------------------------------------------------------------
\usetheme[progressbar=foot, block=fill, sectionpage=progressbar]{metropolis}

% ------------------------------------------------------------------------------
% Core Packages
% ------------------------------------------------------------------------------
\usepackage{tikz}
\usetikzlibrary{shapes,arrows,positioning,calc,decorations.pathreplacing,fit,backgrounds,chains}
\usepackage{booktabs}
\usepackage{array}
\usepackage{graphicx}
\usepackage{xcolor}
\usepackage[sfdefault]{plex-sans}   % IBM Plex Sans — compact, modern, tech-themed
\usepackage[T1]{fontenc}
\usepackage{fontawesome5}          % real vector icons

% Conditional text boxes (more flexible than basic beamer blocks)
\usepackage{tcolorbox}
\tcbuselibrary{skins,breakable}

% ------------------------------------------------------------------------------
% Bibliography (loaded in all modes for citations)
% ------------------------------------------------------------------------------
\usepackage[style=authoryear,backend=biber,natbib=true,maxcitenames=2]{biblatex}
\addbibresource{refs.bib}

% ------------------------------------------------------------------------------
% Article Mode Adjustments
% ------------------------------------------------------------------------------
% When compiling as article (using beamerarticle class), apply these settings
\mode<article>{
    % Use standard article sections instead of beamer frames
    \usepackage{fullpage}
    \usepackage{hyperref}
    \hypersetup{
        colorlinks=true,
        linkcolor=blue,
        citecolor=blue,
        urlcolor=blue
    }

    % Redefine frame environment for article mode
    % Frames become subsections in article mode
    \renewenvironment{frame}[1][]{
        \par\medskip
        \noindent\textbf{#1}\par\medskip
    }{}
}

% Presentation mode specific settings
\mode<presentation>{
    \setlength{\parskip}{0pt}

    % Show a TOC slide at the start of each section, current section highlighted
    \AtBeginSection[]{
        \begin{frame}{Lecture Outline}
            \tableofcontents[currentsection, hideallsubsections]
        \end{frame}
    }

    % FiraSans has ~15% taller metrics than Computer Modern; pull line spacing back
    % to avoid the massive overflows Metropolis + FiraSans produce by default.
    \renewcommand{\baselinestretch}{0.88}

    % Tighter list spacing
    \setbeamertemplate{itemize/enumerate body begin}{\setlength{\itemsep}{2pt}\setlength{\parsep}{0pt}}
    \setbeamertemplate{itemize/enumerate subbody begin}{\setlength{\itemsep}{1pt}\setlength{\parsep}{0pt}}
}

% ------------------------------------------------------------------------------
% Standard Colors for Boxes
% ------------------------------------------------------------------------------
% Consistent color scheme across all lectures for semantic elements
\definecolor{boxblue}{HTML}{1A4A7A}       % Arguments, concepts
\definecolor{boxgreen}{HTML}{2E7D32}      % Pro arguments, supporting points
\definecolor{boxred}{HTML}{C0392B}        % Objections, warnings
\definecolor{boxgold}{HTML}{B7770D}       % Case studies, examples
\definecolor{boxgray}{HTML}{555F6D}       % Quotes, citations
\definecolor{boxpurple}{HTML}{6A3093}     % Secondary accent
% Metropolis structural colors
\definecolor{mPrimary}{HTML}{1E3A5F}      % Deep navy — frametitle / title bg
\definecolor{mAccent}{HTML}{D4880A}       % Warm amber — progress bar / separators

% ------------------------------------------------------------------------------
% Beamer Theme Color Setup
% ------------------------------------------------------------------------------
% Applied once here; individual lectures should NOT override these.
\mode<presentation>{
    % Metropolis palette primary drives frametitle + title frame background
    \setbeamercolor{palette primary}{bg=mPrimary, fg=white}
    \setbeamercolor{frametitle}{bg=mPrimary, fg=white}
    % Title page element colors — must be explicit so they stay white
    % when a \usebackgroundtemplate image replaces the default navy fill
    \setbeamercolor{title}{fg=white}
    \setbeamercolor{subtitle}{fg=mAccent}
    \setbeamercolor{author}{fg=white}
    \setbeamercolor{institute}{fg=white!85!black}
    \setbeamercolor{date}{fg=white!75!black}
    % Progress bar and title separator use accent color
    \setbeamercolor{progress bar}{fg=mAccent, bg=mAccent!25}
    \setbeamercolor{title separator}{fg=mAccent}
    \setbeamercolor{alerted text}{fg=boxred}
    % Standard beamer blocks
    \setbeamercolor{block title}{fg=white, bg=boxblue}
    \setbeamercolor{block body}{bg=boxblue!8}
    \setbeamercolor{block title alerted}{fg=white, bg=boxred}
    \setbeamercolor{block body alerted}{bg=boxred!8}
    \setbeamercolor{block title example}{fg=white, bg=boxgreen}
    \setbeamercolor{block body example}{bg=boxgreen!5}
}

% ------------------------------------------------------------------------------
% Standard Custom Boxes
% ------------------------------------------------------------------------------
% Uniform boxes used across all lectures - do not redefine in individual files

% Argument/Reasoning Box (logical arguments in standard form)
\newtcolorbox{argumentbox}[1][Argument]{
    colback=boxblue!5,
    colframe=boxblue,
    fonttitle=\bfseries,
    title={#1},
    boxrule=1pt,
    arc=2pt
}

% Quotation Box (for philosophical quotes, sources)
\newtcolorbox{quotebox}[1][Quote]{
    colback=gray!10,
    colframe=boxgray,
    fonttitle=\itshape,
    title={#1},
    boxrule=0.5pt,
    arc=2pt,
    left=10pt,
    right=10pt,
    leftrule=3pt
}

% Case Study/Example Box
\newtcolorbox{casestudybox}[1][Case Study]{
    colback=boxgold!10,
    colframe=boxgold,
    fonttitle=\bfseries,
    title={#1},
    boxrule=1pt,
    arc=2pt
}

% Objection/Counterargument Box
\newtcolorbox{objectionbox}[1][Objection]{
    colback=boxred!5,
    colframe=boxred,
    fonttitle=\bfseries,
    title={#1},
    boxrule=1pt,
    arc=2pt
}

% Pro/Supporting Argument Box
\newtcolorbox{probox}[1][Supporting Argument]{
    colback=boxgreen!5,
    colframe=boxgreen,
    fonttitle=\bfseries,
    title={#1},
    boxrule=1pt,
    arc=2pt
}

% Concept/Definition Box
\newtcolorbox{conceptbox}[1][Key Concept]{
    colback=boxblue!5,
    colframe=boxblue,
    fonttitle=\bfseries,
    title={#1},
    boxrule=1pt,
    arc=2pt
}

% Discussion Box (for group activities, questions)
\newtcolorbox{discussionbox}[1][Discussion]{
    colback=boxgold!5,
    colframe=boxgold,
    fonttitle=\bfseries,
    title={#1},
    boxrule=1pt,
    arc=2pt
}

% Figure Box (for diagrams with explanations)
\newtcolorbox{figurebox}[1][Figure]{
    colback=boxblue!3,
    colframe=boxblue!70,
    fonttitle=\bfseries,
    title={#1},
    boxrule=0.75pt,
    arc=2pt
}

% ------------------------------------------------------------------------------
% Utility Commands and Icons
% ------------------------------------------------------------------------------

% fontawesome5 icon shorthands (\faXxx commands come from the fontawesome5 package)
\newcommand{\iconQuote}{\faQuoteLeft}
\newcommand{\iconBalance}{\faBalanceScale}
\newcommand{\iconDiscussion}{\faComments}
\newcommand{\iconIdea}{\faLightbulb}
\newcommand{\iconPro}{\faThumbsUp}
\newcommand{\iconCon}{\faThumbsDown}

% Discussion question command (for interactive moments)
\newcommand{\discussionquestion}[1]{%
    \begin{alertblock}{Discussion Question}
    #1
    \end{alertblock}
}

% fontawesome5 defines \faQuoteLeft, \faBalanceScale, etc. natively --- no aliases needed.
\newcommand{\discussion}[1]{\discussionquestion{#1}}

% Simple divider for slide sections
\newcommand{\sectiondivider}{%
    \begin{center}
    \rule{0.8\textwidth}{0.4pt}
    \end{center}
}

% Argument environment (for standard form arguments)
\newenvironment{argument}{%
  \begin{argumentbox}[Argument in Standard Form]
  \begin{enumerate}
}{%
  \end{enumerate}
  \end{argumentbox}
}

% Quote slide environment (combines icon and quotebox)
\newenvironment{quoteslide}[2]{%
  \begin{quotebox}[#1]
  \small
  \textit{``#2''}
}{%
  \end{quotebox}
}

% ==============================================================================
% End of Common Preamble
% ==============================================================================

% 3. Define lecture-specific colors/customizations
% 4. Set title information
% ==============================================================================

% ------------------------------------------------------------------------------
% Theme and Basic Setup
% ------------------------------------------------------------------------------
\usetheme[progressbar=foot, block=fill, sectionpage=progressbar]{metropolis}

% ------------------------------------------------------------------------------
% Core Packages
% ------------------------------------------------------------------------------
\usepackage{tikz}
\usetikzlibrary{shapes,arrows,positioning,calc,decorations.pathreplacing,fit,backgrounds,chains}
\usepackage{booktabs}
\usepackage{array}
\usepackage{graphicx}
\usepackage{xcolor}
\usepackage[sfdefault]{plex-sans}   % IBM Plex Sans — compact, modern, tech-themed
\usepackage[T1]{fontenc}
\usepackage{fontawesome5}          % real vector icons

% Conditional text boxes (more flexible than basic beamer blocks)
\usepackage{tcolorbox}
\tcbuselibrary{skins,breakable}

% ------------------------------------------------------------------------------
% Bibliography (loaded in all modes for citations)
% ------------------------------------------------------------------------------
\usepackage[style=authoryear,backend=biber,natbib=true,maxcitenames=2]{biblatex}
\addbibresource{refs.bib}

% ------------------------------------------------------------------------------
% Article Mode Adjustments
% ------------------------------------------------------------------------------
% When compiling as article (using beamerarticle class), apply these settings
\mode<article>{
    % Use standard article sections instead of beamer frames
    \usepackage{fullpage}
    \usepackage{hyperref}
    \hypersetup{
        colorlinks=true,
        linkcolor=blue,
        citecolor=blue,
        urlcolor=blue
    }

    % Redefine frame environment for article mode
    % Frames become subsections in article mode
    \renewenvironment{frame}[1][]{
        \par\medskip
        \noindent\textbf{#1}\par\medskip
    }{}
}

% Presentation mode specific settings
\mode<presentation>{
    \setlength{\parskip}{0pt}

    % Show a TOC slide at the start of each section, current section highlighted
    \AtBeginSection[]{
        \begin{frame}{Lecture Outline}
            \tableofcontents[currentsection, hideallsubsections]
        \end{frame}
    }

    % FiraSans has ~15% taller metrics than Computer Modern; pull line spacing back
    % to avoid the massive overflows Metropolis + FiraSans produce by default.
    \renewcommand{\baselinestretch}{0.88}

    % Tighter list spacing
    \setbeamertemplate{itemize/enumerate body begin}{\setlength{\itemsep}{2pt}\setlength{\parsep}{0pt}}
    \setbeamertemplate{itemize/enumerate subbody begin}{\setlength{\itemsep}{1pt}\setlength{\parsep}{0pt}}
}

% ------------------------------------------------------------------------------
% Standard Colors for Boxes
% ------------------------------------------------------------------------------
% Consistent color scheme across all lectures for semantic elements
\definecolor{boxblue}{HTML}{1A4A7A}       % Arguments, concepts
\definecolor{boxgreen}{HTML}{2E7D32}      % Pro arguments, supporting points
\definecolor{boxred}{HTML}{C0392B}        % Objections, warnings
\definecolor{boxgold}{HTML}{B7770D}       % Case studies, examples
\definecolor{boxgray}{HTML}{555F6D}       % Quotes, citations
\definecolor{boxpurple}{HTML}{6A3093}     % Secondary accent
% Metropolis structural colors
\definecolor{mPrimary}{HTML}{1E3A5F}      % Deep navy — frametitle / title bg
\definecolor{mAccent}{HTML}{D4880A}       % Warm amber — progress bar / separators

% ------------------------------------------------------------------------------
% Beamer Theme Color Setup
% ------------------------------------------------------------------------------
% Applied once here; individual lectures should NOT override these.
\mode<presentation>{
    % Metropolis palette primary drives frametitle + title frame background
    \setbeamercolor{palette primary}{bg=mPrimary, fg=white}
    \setbeamercolor{frametitle}{bg=mPrimary, fg=white}
    % Title page element colors — must be explicit so they stay white
    % when a \usebackgroundtemplate image replaces the default navy fill
    \setbeamercolor{title}{fg=white}
    \setbeamercolor{subtitle}{fg=mAccent}
    \setbeamercolor{author}{fg=white}
    \setbeamercolor{institute}{fg=white!85!black}
    \setbeamercolor{date}{fg=white!75!black}
    % Progress bar and title separator use accent color
    \setbeamercolor{progress bar}{fg=mAccent, bg=mAccent!25}
    \setbeamercolor{title separator}{fg=mAccent}
    \setbeamercolor{alerted text}{fg=boxred}
    % Standard beamer blocks
    \setbeamercolor{block title}{fg=white, bg=boxblue}
    \setbeamercolor{block body}{bg=boxblue!8}
    \setbeamercolor{block title alerted}{fg=white, bg=boxred}
    \setbeamercolor{block body alerted}{bg=boxred!8}
    \setbeamercolor{block title example}{fg=white, bg=boxgreen}
    \setbeamercolor{block body example}{bg=boxgreen!5}
}

% ------------------------------------------------------------------------------
% Standard Custom Boxes
% ------------------------------------------------------------------------------
% Uniform boxes used across all lectures - do not redefine in individual files

% Argument/Reasoning Box (logical arguments in standard form)
\newtcolorbox{argumentbox}[1][Argument]{
    colback=boxblue!5,
    colframe=boxblue,
    fonttitle=\bfseries,
    title={#1},
    boxrule=1pt,
    arc=2pt
}

% Quotation Box (for philosophical quotes, sources)
\newtcolorbox{quotebox}[1][Quote]{
    colback=gray!10,
    colframe=boxgray,
    fonttitle=\itshape,
    title={#1},
    boxrule=0.5pt,
    arc=2pt,
    left=10pt,
    right=10pt,
    leftrule=3pt
}

% Case Study/Example Box
\newtcolorbox{casestudybox}[1][Case Study]{
    colback=boxgold!10,
    colframe=boxgold,
    fonttitle=\bfseries,
    title={#1},
    boxrule=1pt,
    arc=2pt
}

% Objection/Counterargument Box
\newtcolorbox{objectionbox}[1][Objection]{
    colback=boxred!5,
    colframe=boxred,
    fonttitle=\bfseries,
    title={#1},
    boxrule=1pt,
    arc=2pt
}

% Pro/Supporting Argument Box
\newtcolorbox{probox}[1][Supporting Argument]{
    colback=boxgreen!5,
    colframe=boxgreen,
    fonttitle=\bfseries,
    title={#1},
    boxrule=1pt,
    arc=2pt
}

% Concept/Definition Box
\newtcolorbox{conceptbox}[1][Key Concept]{
    colback=boxblue!5,
    colframe=boxblue,
    fonttitle=\bfseries,
    title={#1},
    boxrule=1pt,
    arc=2pt
}

% Discussion Box (for group activities, questions)
\newtcolorbox{discussionbox}[1][Discussion]{
    colback=boxgold!5,
    colframe=boxgold,
    fonttitle=\bfseries,
    title={#1},
    boxrule=1pt,
    arc=2pt
}

% Figure Box (for diagrams with explanations)
\newtcolorbox{figurebox}[1][Figure]{
    colback=boxblue!3,
    colframe=boxblue!70,
    fonttitle=\bfseries,
    title={#1},
    boxrule=0.75pt,
    arc=2pt
}

% ------------------------------------------------------------------------------
% Utility Commands and Icons
% ------------------------------------------------------------------------------

% fontawesome5 icon shorthands (\faXxx commands come from the fontawesome5 package)
\newcommand{\iconQuote}{\faQuoteLeft}
\newcommand{\iconBalance}{\faBalanceScale}
\newcommand{\iconDiscussion}{\faComments}
\newcommand{\iconIdea}{\faLightbulb}
\newcommand{\iconPro}{\faThumbsUp}
\newcommand{\iconCon}{\faThumbsDown}

% Discussion question command (for interactive moments)
\newcommand{\discussionquestion}[1]{%
    \begin{alertblock}{Discussion Question}
    #1
    \end{alertblock}
}

% fontawesome5 defines \faQuoteLeft, \faBalanceScale, etc. natively --- no aliases needed.
\newcommand{\discussion}[1]{\discussionquestion{#1}}

% Simple divider for slide sections
\newcommand{\sectiondivider}{%
    \begin{center}
    \rule{0.8\textwidth}{0.4pt}
    \end{center}
}

% Argument environment (for standard form arguments)
\newenvironment{argument}{%
  \begin{argumentbox}[Argument in Standard Form]
  \begin{enumerate}
}{%
  \end{enumerate}
  \end{argumentbox}
}

% Quote slide environment (combines icon and quotebox)
\newenvironment{quoteslide}[2]{%
  \begin{quotebox}[#1]
  \small
  \textit{``#2''}
}{%
  \end{quotebox}
}

% ==============================================================================
% End of Common Preamble
% ==============================================================================

%
% Individual lecture files should:
% 1. Declare \documentclass (beamer or beamerarticle)
% 2. % ==============================================================================
% Common Preamble for IT Ethics Lectures
% ==============================================================================
% This file contains shared configuration for all lecture files.
% Include this after \documentclass using: % ==============================================================================
% Common Preamble for IT Ethics Lectures
% ==============================================================================
% This file contains shared configuration for all lecture files.
% Include this after \documentclass using: % ==============================================================================
% Common Preamble for IT Ethics Lectures
% ==============================================================================
% This file contains shared configuration for all lecture files.
% Include this after \documentclass using: \input{lecture_preamble.tex}
%
% Individual lecture files should:
% 1. Declare \documentclass (beamer or beamerarticle)
% 2. \input{lecture_preamble.tex}
% 3. Define lecture-specific colors/customizations
% 4. Set title information
% ==============================================================================

% ------------------------------------------------------------------------------
% Theme and Basic Setup
% ------------------------------------------------------------------------------
\usetheme[progressbar=foot, block=fill, sectionpage=progressbar]{metropolis}

% ------------------------------------------------------------------------------
% Core Packages
% ------------------------------------------------------------------------------
\usepackage{tikz}
\usetikzlibrary{shapes,arrows,positioning,calc,decorations.pathreplacing,fit,backgrounds,chains}
\usepackage{booktabs}
\usepackage{array}
\usepackage{graphicx}
\usepackage{xcolor}
\usepackage[sfdefault]{plex-sans}   % IBM Plex Sans — compact, modern, tech-themed
\usepackage[T1]{fontenc}
\usepackage{fontawesome5}          % real vector icons

% Conditional text boxes (more flexible than basic beamer blocks)
\usepackage{tcolorbox}
\tcbuselibrary{skins,breakable}

% ------------------------------------------------------------------------------
% Bibliography (loaded in all modes for citations)
% ------------------------------------------------------------------------------
\usepackage[style=authoryear,backend=biber,natbib=true,maxcitenames=2]{biblatex}
\addbibresource{refs.bib}

% ------------------------------------------------------------------------------
% Article Mode Adjustments
% ------------------------------------------------------------------------------
% When compiling as article (using beamerarticle class), apply these settings
\mode<article>{
    % Use standard article sections instead of beamer frames
    \usepackage{fullpage}
    \usepackage{hyperref}
    \hypersetup{
        colorlinks=true,
        linkcolor=blue,
        citecolor=blue,
        urlcolor=blue
    }

    % Redefine frame environment for article mode
    % Frames become subsections in article mode
    \renewenvironment{frame}[1][]{
        \par\medskip
        \noindent\textbf{#1}\par\medskip
    }{}
}

% Presentation mode specific settings
\mode<presentation>{
    \setlength{\parskip}{0pt}

    % Show a TOC slide at the start of each section, current section highlighted
    \AtBeginSection[]{
        \begin{frame}{Lecture Outline}
            \tableofcontents[currentsection, hideallsubsections]
        \end{frame}
    }

    % FiraSans has ~15% taller metrics than Computer Modern; pull line spacing back
    % to avoid the massive overflows Metropolis + FiraSans produce by default.
    \renewcommand{\baselinestretch}{0.88}

    % Tighter list spacing
    \setbeamertemplate{itemize/enumerate body begin}{\setlength{\itemsep}{2pt}\setlength{\parsep}{0pt}}
    \setbeamertemplate{itemize/enumerate subbody begin}{\setlength{\itemsep}{1pt}\setlength{\parsep}{0pt}}
}

% ------------------------------------------------------------------------------
% Standard Colors for Boxes
% ------------------------------------------------------------------------------
% Consistent color scheme across all lectures for semantic elements
\definecolor{boxblue}{HTML}{1A4A7A}       % Arguments, concepts
\definecolor{boxgreen}{HTML}{2E7D32}      % Pro arguments, supporting points
\definecolor{boxred}{HTML}{C0392B}        % Objections, warnings
\definecolor{boxgold}{HTML}{B7770D}       % Case studies, examples
\definecolor{boxgray}{HTML}{555F6D}       % Quotes, citations
\definecolor{boxpurple}{HTML}{6A3093}     % Secondary accent
% Metropolis structural colors
\definecolor{mPrimary}{HTML}{1E3A5F}      % Deep navy — frametitle / title bg
\definecolor{mAccent}{HTML}{D4880A}       % Warm amber — progress bar / separators

% ------------------------------------------------------------------------------
% Beamer Theme Color Setup
% ------------------------------------------------------------------------------
% Applied once here; individual lectures should NOT override these.
\mode<presentation>{
    % Metropolis palette primary drives frametitle + title frame background
    \setbeamercolor{palette primary}{bg=mPrimary, fg=white}
    \setbeamercolor{frametitle}{bg=mPrimary, fg=white}
    % Title page element colors — must be explicit so they stay white
    % when a \usebackgroundtemplate image replaces the default navy fill
    \setbeamercolor{title}{fg=white}
    \setbeamercolor{subtitle}{fg=mAccent}
    \setbeamercolor{author}{fg=white}
    \setbeamercolor{institute}{fg=white!85!black}
    \setbeamercolor{date}{fg=white!75!black}
    % Progress bar and title separator use accent color
    \setbeamercolor{progress bar}{fg=mAccent, bg=mAccent!25}
    \setbeamercolor{title separator}{fg=mAccent}
    \setbeamercolor{alerted text}{fg=boxred}
    % Standard beamer blocks
    \setbeamercolor{block title}{fg=white, bg=boxblue}
    \setbeamercolor{block body}{bg=boxblue!8}
    \setbeamercolor{block title alerted}{fg=white, bg=boxred}
    \setbeamercolor{block body alerted}{bg=boxred!8}
    \setbeamercolor{block title example}{fg=white, bg=boxgreen}
    \setbeamercolor{block body example}{bg=boxgreen!5}
}

% ------------------------------------------------------------------------------
% Standard Custom Boxes
% ------------------------------------------------------------------------------
% Uniform boxes used across all lectures - do not redefine in individual files

% Argument/Reasoning Box (logical arguments in standard form)
\newtcolorbox{argumentbox}[1][Argument]{
    colback=boxblue!5,
    colframe=boxblue,
    fonttitle=\bfseries,
    title={#1},
    boxrule=1pt,
    arc=2pt
}

% Quotation Box (for philosophical quotes, sources)
\newtcolorbox{quotebox}[1][Quote]{
    colback=gray!10,
    colframe=boxgray,
    fonttitle=\itshape,
    title={#1},
    boxrule=0.5pt,
    arc=2pt,
    left=10pt,
    right=10pt,
    leftrule=3pt
}

% Case Study/Example Box
\newtcolorbox{casestudybox}[1][Case Study]{
    colback=boxgold!10,
    colframe=boxgold,
    fonttitle=\bfseries,
    title={#1},
    boxrule=1pt,
    arc=2pt
}

% Objection/Counterargument Box
\newtcolorbox{objectionbox}[1][Objection]{
    colback=boxred!5,
    colframe=boxred,
    fonttitle=\bfseries,
    title={#1},
    boxrule=1pt,
    arc=2pt
}

% Pro/Supporting Argument Box
\newtcolorbox{probox}[1][Supporting Argument]{
    colback=boxgreen!5,
    colframe=boxgreen,
    fonttitle=\bfseries,
    title={#1},
    boxrule=1pt,
    arc=2pt
}

% Concept/Definition Box
\newtcolorbox{conceptbox}[1][Key Concept]{
    colback=boxblue!5,
    colframe=boxblue,
    fonttitle=\bfseries,
    title={#1},
    boxrule=1pt,
    arc=2pt
}

% Discussion Box (for group activities, questions)
\newtcolorbox{discussionbox}[1][Discussion]{
    colback=boxgold!5,
    colframe=boxgold,
    fonttitle=\bfseries,
    title={#1},
    boxrule=1pt,
    arc=2pt
}

% Figure Box (for diagrams with explanations)
\newtcolorbox{figurebox}[1][Figure]{
    colback=boxblue!3,
    colframe=boxblue!70,
    fonttitle=\bfseries,
    title={#1},
    boxrule=0.75pt,
    arc=2pt
}

% ------------------------------------------------------------------------------
% Utility Commands and Icons
% ------------------------------------------------------------------------------

% fontawesome5 icon shorthands (\faXxx commands come from the fontawesome5 package)
\newcommand{\iconQuote}{\faQuoteLeft}
\newcommand{\iconBalance}{\faBalanceScale}
\newcommand{\iconDiscussion}{\faComments}
\newcommand{\iconIdea}{\faLightbulb}
\newcommand{\iconPro}{\faThumbsUp}
\newcommand{\iconCon}{\faThumbsDown}

% Discussion question command (for interactive moments)
\newcommand{\discussionquestion}[1]{%
    \begin{alertblock}{Discussion Question}
    #1
    \end{alertblock}
}

% fontawesome5 defines \faQuoteLeft, \faBalanceScale, etc. natively --- no aliases needed.
\newcommand{\discussion}[1]{\discussionquestion{#1}}

% Simple divider for slide sections
\newcommand{\sectiondivider}{%
    \begin{center}
    \rule{0.8\textwidth}{0.4pt}
    \end{center}
}

% Argument environment (for standard form arguments)
\newenvironment{argument}{%
  \begin{argumentbox}[Argument in Standard Form]
  \begin{enumerate}
}{%
  \end{enumerate}
  \end{argumentbox}
}

% Quote slide environment (combines icon and quotebox)
\newenvironment{quoteslide}[2]{%
  \begin{quotebox}[#1]
  \small
  \textit{``#2''}
}{%
  \end{quotebox}
}

% ==============================================================================
% End of Common Preamble
% ==============================================================================

%
% Individual lecture files should:
% 1. Declare \documentclass (beamer or beamerarticle)
% 2. % ==============================================================================
% Common Preamble for IT Ethics Lectures
% ==============================================================================
% This file contains shared configuration for all lecture files.
% Include this after \documentclass using: \input{lecture_preamble.tex}
%
% Individual lecture files should:
% 1. Declare \documentclass (beamer or beamerarticle)
% 2. \input{lecture_preamble.tex}
% 3. Define lecture-specific colors/customizations
% 4. Set title information
% ==============================================================================

% ------------------------------------------------------------------------------
% Theme and Basic Setup
% ------------------------------------------------------------------------------
\usetheme[progressbar=foot, block=fill, sectionpage=progressbar]{metropolis}

% ------------------------------------------------------------------------------
% Core Packages
% ------------------------------------------------------------------------------
\usepackage{tikz}
\usetikzlibrary{shapes,arrows,positioning,calc,decorations.pathreplacing,fit,backgrounds,chains}
\usepackage{booktabs}
\usepackage{array}
\usepackage{graphicx}
\usepackage{xcolor}
\usepackage[sfdefault]{plex-sans}   % IBM Plex Sans — compact, modern, tech-themed
\usepackage[T1]{fontenc}
\usepackage{fontawesome5}          % real vector icons

% Conditional text boxes (more flexible than basic beamer blocks)
\usepackage{tcolorbox}
\tcbuselibrary{skins,breakable}

% ------------------------------------------------------------------------------
% Bibliography (loaded in all modes for citations)
% ------------------------------------------------------------------------------
\usepackage[style=authoryear,backend=biber,natbib=true,maxcitenames=2]{biblatex}
\addbibresource{refs.bib}

% ------------------------------------------------------------------------------
% Article Mode Adjustments
% ------------------------------------------------------------------------------
% When compiling as article (using beamerarticle class), apply these settings
\mode<article>{
    % Use standard article sections instead of beamer frames
    \usepackage{fullpage}
    \usepackage{hyperref}
    \hypersetup{
        colorlinks=true,
        linkcolor=blue,
        citecolor=blue,
        urlcolor=blue
    }

    % Redefine frame environment for article mode
    % Frames become subsections in article mode
    \renewenvironment{frame}[1][]{
        \par\medskip
        \noindent\textbf{#1}\par\medskip
    }{}
}

% Presentation mode specific settings
\mode<presentation>{
    \setlength{\parskip}{0pt}

    % Show a TOC slide at the start of each section, current section highlighted
    \AtBeginSection[]{
        \begin{frame}{Lecture Outline}
            \tableofcontents[currentsection, hideallsubsections]
        \end{frame}
    }

    % FiraSans has ~15% taller metrics than Computer Modern; pull line spacing back
    % to avoid the massive overflows Metropolis + FiraSans produce by default.
    \renewcommand{\baselinestretch}{0.88}

    % Tighter list spacing
    \setbeamertemplate{itemize/enumerate body begin}{\setlength{\itemsep}{2pt}\setlength{\parsep}{0pt}}
    \setbeamertemplate{itemize/enumerate subbody begin}{\setlength{\itemsep}{1pt}\setlength{\parsep}{0pt}}
}

% ------------------------------------------------------------------------------
% Standard Colors for Boxes
% ------------------------------------------------------------------------------
% Consistent color scheme across all lectures for semantic elements
\definecolor{boxblue}{HTML}{1A4A7A}       % Arguments, concepts
\definecolor{boxgreen}{HTML}{2E7D32}      % Pro arguments, supporting points
\definecolor{boxred}{HTML}{C0392B}        % Objections, warnings
\definecolor{boxgold}{HTML}{B7770D}       % Case studies, examples
\definecolor{boxgray}{HTML}{555F6D}       % Quotes, citations
\definecolor{boxpurple}{HTML}{6A3093}     % Secondary accent
% Metropolis structural colors
\definecolor{mPrimary}{HTML}{1E3A5F}      % Deep navy — frametitle / title bg
\definecolor{mAccent}{HTML}{D4880A}       % Warm amber — progress bar / separators

% ------------------------------------------------------------------------------
% Beamer Theme Color Setup
% ------------------------------------------------------------------------------
% Applied once here; individual lectures should NOT override these.
\mode<presentation>{
    % Metropolis palette primary drives frametitle + title frame background
    \setbeamercolor{palette primary}{bg=mPrimary, fg=white}
    \setbeamercolor{frametitle}{bg=mPrimary, fg=white}
    % Title page element colors — must be explicit so they stay white
    % when a \usebackgroundtemplate image replaces the default navy fill
    \setbeamercolor{title}{fg=white}
    \setbeamercolor{subtitle}{fg=mAccent}
    \setbeamercolor{author}{fg=white}
    \setbeamercolor{institute}{fg=white!85!black}
    \setbeamercolor{date}{fg=white!75!black}
    % Progress bar and title separator use accent color
    \setbeamercolor{progress bar}{fg=mAccent, bg=mAccent!25}
    \setbeamercolor{title separator}{fg=mAccent}
    \setbeamercolor{alerted text}{fg=boxred}
    % Standard beamer blocks
    \setbeamercolor{block title}{fg=white, bg=boxblue}
    \setbeamercolor{block body}{bg=boxblue!8}
    \setbeamercolor{block title alerted}{fg=white, bg=boxred}
    \setbeamercolor{block body alerted}{bg=boxred!8}
    \setbeamercolor{block title example}{fg=white, bg=boxgreen}
    \setbeamercolor{block body example}{bg=boxgreen!5}
}

% ------------------------------------------------------------------------------
% Standard Custom Boxes
% ------------------------------------------------------------------------------
% Uniform boxes used across all lectures - do not redefine in individual files

% Argument/Reasoning Box (logical arguments in standard form)
\newtcolorbox{argumentbox}[1][Argument]{
    colback=boxblue!5,
    colframe=boxblue,
    fonttitle=\bfseries,
    title={#1},
    boxrule=1pt,
    arc=2pt
}

% Quotation Box (for philosophical quotes, sources)
\newtcolorbox{quotebox}[1][Quote]{
    colback=gray!10,
    colframe=boxgray,
    fonttitle=\itshape,
    title={#1},
    boxrule=0.5pt,
    arc=2pt,
    left=10pt,
    right=10pt,
    leftrule=3pt
}

% Case Study/Example Box
\newtcolorbox{casestudybox}[1][Case Study]{
    colback=boxgold!10,
    colframe=boxgold,
    fonttitle=\bfseries,
    title={#1},
    boxrule=1pt,
    arc=2pt
}

% Objection/Counterargument Box
\newtcolorbox{objectionbox}[1][Objection]{
    colback=boxred!5,
    colframe=boxred,
    fonttitle=\bfseries,
    title={#1},
    boxrule=1pt,
    arc=2pt
}

% Pro/Supporting Argument Box
\newtcolorbox{probox}[1][Supporting Argument]{
    colback=boxgreen!5,
    colframe=boxgreen,
    fonttitle=\bfseries,
    title={#1},
    boxrule=1pt,
    arc=2pt
}

% Concept/Definition Box
\newtcolorbox{conceptbox}[1][Key Concept]{
    colback=boxblue!5,
    colframe=boxblue,
    fonttitle=\bfseries,
    title={#1},
    boxrule=1pt,
    arc=2pt
}

% Discussion Box (for group activities, questions)
\newtcolorbox{discussionbox}[1][Discussion]{
    colback=boxgold!5,
    colframe=boxgold,
    fonttitle=\bfseries,
    title={#1},
    boxrule=1pt,
    arc=2pt
}

% Figure Box (for diagrams with explanations)
\newtcolorbox{figurebox}[1][Figure]{
    colback=boxblue!3,
    colframe=boxblue!70,
    fonttitle=\bfseries,
    title={#1},
    boxrule=0.75pt,
    arc=2pt
}

% ------------------------------------------------------------------------------
% Utility Commands and Icons
% ------------------------------------------------------------------------------

% fontawesome5 icon shorthands (\faXxx commands come from the fontawesome5 package)
\newcommand{\iconQuote}{\faQuoteLeft}
\newcommand{\iconBalance}{\faBalanceScale}
\newcommand{\iconDiscussion}{\faComments}
\newcommand{\iconIdea}{\faLightbulb}
\newcommand{\iconPro}{\faThumbsUp}
\newcommand{\iconCon}{\faThumbsDown}

% Discussion question command (for interactive moments)
\newcommand{\discussionquestion}[1]{%
    \begin{alertblock}{Discussion Question}
    #1
    \end{alertblock}
}

% fontawesome5 defines \faQuoteLeft, \faBalanceScale, etc. natively --- no aliases needed.
\newcommand{\discussion}[1]{\discussionquestion{#1}}

% Simple divider for slide sections
\newcommand{\sectiondivider}{%
    \begin{center}
    \rule{0.8\textwidth}{0.4pt}
    \end{center}
}

% Argument environment (for standard form arguments)
\newenvironment{argument}{%
  \begin{argumentbox}[Argument in Standard Form]
  \begin{enumerate}
}{%
  \end{enumerate}
  \end{argumentbox}
}

% Quote slide environment (combines icon and quotebox)
\newenvironment{quoteslide}[2]{%
  \begin{quotebox}[#1]
  \small
  \textit{``#2''}
}{%
  \end{quotebox}
}

% ==============================================================================
% End of Common Preamble
% ==============================================================================

% 3. Define lecture-specific colors/customizations
% 4. Set title information
% ==============================================================================

% ------------------------------------------------------------------------------
% Theme and Basic Setup
% ------------------------------------------------------------------------------
\usetheme[progressbar=foot, block=fill, sectionpage=progressbar]{metropolis}

% ------------------------------------------------------------------------------
% Core Packages
% ------------------------------------------------------------------------------
\usepackage{tikz}
\usetikzlibrary{shapes,arrows,positioning,calc,decorations.pathreplacing,fit,backgrounds,chains}
\usepackage{booktabs}
\usepackage{array}
\usepackage{graphicx}
\usepackage{xcolor}
\usepackage[sfdefault]{plex-sans}   % IBM Plex Sans — compact, modern, tech-themed
\usepackage[T1]{fontenc}
\usepackage{fontawesome5}          % real vector icons

% Conditional text boxes (more flexible than basic beamer blocks)
\usepackage{tcolorbox}
\tcbuselibrary{skins,breakable}

% ------------------------------------------------------------------------------
% Bibliography (loaded in all modes for citations)
% ------------------------------------------------------------------------------
\usepackage[style=authoryear,backend=biber,natbib=true,maxcitenames=2]{biblatex}
\addbibresource{refs.bib}

% ------------------------------------------------------------------------------
% Article Mode Adjustments
% ------------------------------------------------------------------------------
% When compiling as article (using beamerarticle class), apply these settings
\mode<article>{
    % Use standard article sections instead of beamer frames
    \usepackage{fullpage}
    \usepackage{hyperref}
    \hypersetup{
        colorlinks=true,
        linkcolor=blue,
        citecolor=blue,
        urlcolor=blue
    }

    % Redefine frame environment for article mode
    % Frames become subsections in article mode
    \renewenvironment{frame}[1][]{
        \par\medskip
        \noindent\textbf{#1}\par\medskip
    }{}
}

% Presentation mode specific settings
\mode<presentation>{
    \setlength{\parskip}{0pt}

    % Show a TOC slide at the start of each section, current section highlighted
    \AtBeginSection[]{
        \begin{frame}{Lecture Outline}
            \tableofcontents[currentsection, hideallsubsections]
        \end{frame}
    }

    % FiraSans has ~15% taller metrics than Computer Modern; pull line spacing back
    % to avoid the massive overflows Metropolis + FiraSans produce by default.
    \renewcommand{\baselinestretch}{0.88}

    % Tighter list spacing
    \setbeamertemplate{itemize/enumerate body begin}{\setlength{\itemsep}{2pt}\setlength{\parsep}{0pt}}
    \setbeamertemplate{itemize/enumerate subbody begin}{\setlength{\itemsep}{1pt}\setlength{\parsep}{0pt}}
}

% ------------------------------------------------------------------------------
% Standard Colors for Boxes
% ------------------------------------------------------------------------------
% Consistent color scheme across all lectures for semantic elements
\definecolor{boxblue}{HTML}{1A4A7A}       % Arguments, concepts
\definecolor{boxgreen}{HTML}{2E7D32}      % Pro arguments, supporting points
\definecolor{boxred}{HTML}{C0392B}        % Objections, warnings
\definecolor{boxgold}{HTML}{B7770D}       % Case studies, examples
\definecolor{boxgray}{HTML}{555F6D}       % Quotes, citations
\definecolor{boxpurple}{HTML}{6A3093}     % Secondary accent
% Metropolis structural colors
\definecolor{mPrimary}{HTML}{1E3A5F}      % Deep navy — frametitle / title bg
\definecolor{mAccent}{HTML}{D4880A}       % Warm amber — progress bar / separators

% ------------------------------------------------------------------------------
% Beamer Theme Color Setup
% ------------------------------------------------------------------------------
% Applied once here; individual lectures should NOT override these.
\mode<presentation>{
    % Metropolis palette primary drives frametitle + title frame background
    \setbeamercolor{palette primary}{bg=mPrimary, fg=white}
    \setbeamercolor{frametitle}{bg=mPrimary, fg=white}
    % Title page element colors — must be explicit so they stay white
    % when a \usebackgroundtemplate image replaces the default navy fill
    \setbeamercolor{title}{fg=white}
    \setbeamercolor{subtitle}{fg=mAccent}
    \setbeamercolor{author}{fg=white}
    \setbeamercolor{institute}{fg=white!85!black}
    \setbeamercolor{date}{fg=white!75!black}
    % Progress bar and title separator use accent color
    \setbeamercolor{progress bar}{fg=mAccent, bg=mAccent!25}
    \setbeamercolor{title separator}{fg=mAccent}
    \setbeamercolor{alerted text}{fg=boxred}
    % Standard beamer blocks
    \setbeamercolor{block title}{fg=white, bg=boxblue}
    \setbeamercolor{block body}{bg=boxblue!8}
    \setbeamercolor{block title alerted}{fg=white, bg=boxred}
    \setbeamercolor{block body alerted}{bg=boxred!8}
    \setbeamercolor{block title example}{fg=white, bg=boxgreen}
    \setbeamercolor{block body example}{bg=boxgreen!5}
}

% ------------------------------------------------------------------------------
% Standard Custom Boxes
% ------------------------------------------------------------------------------
% Uniform boxes used across all lectures - do not redefine in individual files

% Argument/Reasoning Box (logical arguments in standard form)
\newtcolorbox{argumentbox}[1][Argument]{
    colback=boxblue!5,
    colframe=boxblue,
    fonttitle=\bfseries,
    title={#1},
    boxrule=1pt,
    arc=2pt
}

% Quotation Box (for philosophical quotes, sources)
\newtcolorbox{quotebox}[1][Quote]{
    colback=gray!10,
    colframe=boxgray,
    fonttitle=\itshape,
    title={#1},
    boxrule=0.5pt,
    arc=2pt,
    left=10pt,
    right=10pt,
    leftrule=3pt
}

% Case Study/Example Box
\newtcolorbox{casestudybox}[1][Case Study]{
    colback=boxgold!10,
    colframe=boxgold,
    fonttitle=\bfseries,
    title={#1},
    boxrule=1pt,
    arc=2pt
}

% Objection/Counterargument Box
\newtcolorbox{objectionbox}[1][Objection]{
    colback=boxred!5,
    colframe=boxred,
    fonttitle=\bfseries,
    title={#1},
    boxrule=1pt,
    arc=2pt
}

% Pro/Supporting Argument Box
\newtcolorbox{probox}[1][Supporting Argument]{
    colback=boxgreen!5,
    colframe=boxgreen,
    fonttitle=\bfseries,
    title={#1},
    boxrule=1pt,
    arc=2pt
}

% Concept/Definition Box
\newtcolorbox{conceptbox}[1][Key Concept]{
    colback=boxblue!5,
    colframe=boxblue,
    fonttitle=\bfseries,
    title={#1},
    boxrule=1pt,
    arc=2pt
}

% Discussion Box (for group activities, questions)
\newtcolorbox{discussionbox}[1][Discussion]{
    colback=boxgold!5,
    colframe=boxgold,
    fonttitle=\bfseries,
    title={#1},
    boxrule=1pt,
    arc=2pt
}

% Figure Box (for diagrams with explanations)
\newtcolorbox{figurebox}[1][Figure]{
    colback=boxblue!3,
    colframe=boxblue!70,
    fonttitle=\bfseries,
    title={#1},
    boxrule=0.75pt,
    arc=2pt
}

% ------------------------------------------------------------------------------
% Utility Commands and Icons
% ------------------------------------------------------------------------------

% fontawesome5 icon shorthands (\faXxx commands come from the fontawesome5 package)
\newcommand{\iconQuote}{\faQuoteLeft}
\newcommand{\iconBalance}{\faBalanceScale}
\newcommand{\iconDiscussion}{\faComments}
\newcommand{\iconIdea}{\faLightbulb}
\newcommand{\iconPro}{\faThumbsUp}
\newcommand{\iconCon}{\faThumbsDown}

% Discussion question command (for interactive moments)
\newcommand{\discussionquestion}[1]{%
    \begin{alertblock}{Discussion Question}
    #1
    \end{alertblock}
}

% fontawesome5 defines \faQuoteLeft, \faBalanceScale, etc. natively --- no aliases needed.
\newcommand{\discussion}[1]{\discussionquestion{#1}}

% Simple divider for slide sections
\newcommand{\sectiondivider}{%
    \begin{center}
    \rule{0.8\textwidth}{0.4pt}
    \end{center}
}

% Argument environment (for standard form arguments)
\newenvironment{argument}{%
  \begin{argumentbox}[Argument in Standard Form]
  \begin{enumerate}
}{%
  \end{enumerate}
  \end{argumentbox}
}

% Quote slide environment (combines icon and quotebox)
\newenvironment{quoteslide}[2]{%
  \begin{quotebox}[#1]
  \small
  \textit{``#2''}
}{%
  \end{quotebox}
}

% ==============================================================================
% End of Common Preamble
% ==============================================================================

%
% Individual lecture files should:
% 1. Declare \documentclass (beamer or beamerarticle)
% 2. % ==============================================================================
% Common Preamble for IT Ethics Lectures
% ==============================================================================
% This file contains shared configuration for all lecture files.
% Include this after \documentclass using: % ==============================================================================
% Common Preamble for IT Ethics Lectures
% ==============================================================================
% This file contains shared configuration for all lecture files.
% Include this after \documentclass using: \input{lecture_preamble.tex}
%
% Individual lecture files should:
% 1. Declare \documentclass (beamer or beamerarticle)
% 2. \input{lecture_preamble.tex}
% 3. Define lecture-specific colors/customizations
% 4. Set title information
% ==============================================================================

% ------------------------------------------------------------------------------
% Theme and Basic Setup
% ------------------------------------------------------------------------------
\usetheme[progressbar=foot, block=fill, sectionpage=progressbar]{metropolis}

% ------------------------------------------------------------------------------
% Core Packages
% ------------------------------------------------------------------------------
\usepackage{tikz}
\usetikzlibrary{shapes,arrows,positioning,calc,decorations.pathreplacing,fit,backgrounds,chains}
\usepackage{booktabs}
\usepackage{array}
\usepackage{graphicx}
\usepackage{xcolor}
\usepackage[sfdefault]{plex-sans}   % IBM Plex Sans — compact, modern, tech-themed
\usepackage[T1]{fontenc}
\usepackage{fontawesome5}          % real vector icons

% Conditional text boxes (more flexible than basic beamer blocks)
\usepackage{tcolorbox}
\tcbuselibrary{skins,breakable}

% ------------------------------------------------------------------------------
% Bibliography (loaded in all modes for citations)
% ------------------------------------------------------------------------------
\usepackage[style=authoryear,backend=biber,natbib=true,maxcitenames=2]{biblatex}
\addbibresource{refs.bib}

% ------------------------------------------------------------------------------
% Article Mode Adjustments
% ------------------------------------------------------------------------------
% When compiling as article (using beamerarticle class), apply these settings
\mode<article>{
    % Use standard article sections instead of beamer frames
    \usepackage{fullpage}
    \usepackage{hyperref}
    \hypersetup{
        colorlinks=true,
        linkcolor=blue,
        citecolor=blue,
        urlcolor=blue
    }

    % Redefine frame environment for article mode
    % Frames become subsections in article mode
    \renewenvironment{frame}[1][]{
        \par\medskip
        \noindent\textbf{#1}\par\medskip
    }{}
}

% Presentation mode specific settings
\mode<presentation>{
    \setlength{\parskip}{0pt}

    % Show a TOC slide at the start of each section, current section highlighted
    \AtBeginSection[]{
        \begin{frame}{Lecture Outline}
            \tableofcontents[currentsection, hideallsubsections]
        \end{frame}
    }

    % FiraSans has ~15% taller metrics than Computer Modern; pull line spacing back
    % to avoid the massive overflows Metropolis + FiraSans produce by default.
    \renewcommand{\baselinestretch}{0.88}

    % Tighter list spacing
    \setbeamertemplate{itemize/enumerate body begin}{\setlength{\itemsep}{2pt}\setlength{\parsep}{0pt}}
    \setbeamertemplate{itemize/enumerate subbody begin}{\setlength{\itemsep}{1pt}\setlength{\parsep}{0pt}}
}

% ------------------------------------------------------------------------------
% Standard Colors for Boxes
% ------------------------------------------------------------------------------
% Consistent color scheme across all lectures for semantic elements
\definecolor{boxblue}{HTML}{1A4A7A}       % Arguments, concepts
\definecolor{boxgreen}{HTML}{2E7D32}      % Pro arguments, supporting points
\definecolor{boxred}{HTML}{C0392B}        % Objections, warnings
\definecolor{boxgold}{HTML}{B7770D}       % Case studies, examples
\definecolor{boxgray}{HTML}{555F6D}       % Quotes, citations
\definecolor{boxpurple}{HTML}{6A3093}     % Secondary accent
% Metropolis structural colors
\definecolor{mPrimary}{HTML}{1E3A5F}      % Deep navy — frametitle / title bg
\definecolor{mAccent}{HTML}{D4880A}       % Warm amber — progress bar / separators

% ------------------------------------------------------------------------------
% Beamer Theme Color Setup
% ------------------------------------------------------------------------------
% Applied once here; individual lectures should NOT override these.
\mode<presentation>{
    % Metropolis palette primary drives frametitle + title frame background
    \setbeamercolor{palette primary}{bg=mPrimary, fg=white}
    \setbeamercolor{frametitle}{bg=mPrimary, fg=white}
    % Title page element colors — must be explicit so they stay white
    % when a \usebackgroundtemplate image replaces the default navy fill
    \setbeamercolor{title}{fg=white}
    \setbeamercolor{subtitle}{fg=mAccent}
    \setbeamercolor{author}{fg=white}
    \setbeamercolor{institute}{fg=white!85!black}
    \setbeamercolor{date}{fg=white!75!black}
    % Progress bar and title separator use accent color
    \setbeamercolor{progress bar}{fg=mAccent, bg=mAccent!25}
    \setbeamercolor{title separator}{fg=mAccent}
    \setbeamercolor{alerted text}{fg=boxred}
    % Standard beamer blocks
    \setbeamercolor{block title}{fg=white, bg=boxblue}
    \setbeamercolor{block body}{bg=boxblue!8}
    \setbeamercolor{block title alerted}{fg=white, bg=boxred}
    \setbeamercolor{block body alerted}{bg=boxred!8}
    \setbeamercolor{block title example}{fg=white, bg=boxgreen}
    \setbeamercolor{block body example}{bg=boxgreen!5}
}

% ------------------------------------------------------------------------------
% Standard Custom Boxes
% ------------------------------------------------------------------------------
% Uniform boxes used across all lectures - do not redefine in individual files

% Argument/Reasoning Box (logical arguments in standard form)
\newtcolorbox{argumentbox}[1][Argument]{
    colback=boxblue!5,
    colframe=boxblue,
    fonttitle=\bfseries,
    title={#1},
    boxrule=1pt,
    arc=2pt
}

% Quotation Box (for philosophical quotes, sources)
\newtcolorbox{quotebox}[1][Quote]{
    colback=gray!10,
    colframe=boxgray,
    fonttitle=\itshape,
    title={#1},
    boxrule=0.5pt,
    arc=2pt,
    left=10pt,
    right=10pt,
    leftrule=3pt
}

% Case Study/Example Box
\newtcolorbox{casestudybox}[1][Case Study]{
    colback=boxgold!10,
    colframe=boxgold,
    fonttitle=\bfseries,
    title={#1},
    boxrule=1pt,
    arc=2pt
}

% Objection/Counterargument Box
\newtcolorbox{objectionbox}[1][Objection]{
    colback=boxred!5,
    colframe=boxred,
    fonttitle=\bfseries,
    title={#1},
    boxrule=1pt,
    arc=2pt
}

% Pro/Supporting Argument Box
\newtcolorbox{probox}[1][Supporting Argument]{
    colback=boxgreen!5,
    colframe=boxgreen,
    fonttitle=\bfseries,
    title={#1},
    boxrule=1pt,
    arc=2pt
}

% Concept/Definition Box
\newtcolorbox{conceptbox}[1][Key Concept]{
    colback=boxblue!5,
    colframe=boxblue,
    fonttitle=\bfseries,
    title={#1},
    boxrule=1pt,
    arc=2pt
}

% Discussion Box (for group activities, questions)
\newtcolorbox{discussionbox}[1][Discussion]{
    colback=boxgold!5,
    colframe=boxgold,
    fonttitle=\bfseries,
    title={#1},
    boxrule=1pt,
    arc=2pt
}

% Figure Box (for diagrams with explanations)
\newtcolorbox{figurebox}[1][Figure]{
    colback=boxblue!3,
    colframe=boxblue!70,
    fonttitle=\bfseries,
    title={#1},
    boxrule=0.75pt,
    arc=2pt
}

% ------------------------------------------------------------------------------
% Utility Commands and Icons
% ------------------------------------------------------------------------------

% fontawesome5 icon shorthands (\faXxx commands come from the fontawesome5 package)
\newcommand{\iconQuote}{\faQuoteLeft}
\newcommand{\iconBalance}{\faBalanceScale}
\newcommand{\iconDiscussion}{\faComments}
\newcommand{\iconIdea}{\faLightbulb}
\newcommand{\iconPro}{\faThumbsUp}
\newcommand{\iconCon}{\faThumbsDown}

% Discussion question command (for interactive moments)
\newcommand{\discussionquestion}[1]{%
    \begin{alertblock}{Discussion Question}
    #1
    \end{alertblock}
}

% fontawesome5 defines \faQuoteLeft, \faBalanceScale, etc. natively --- no aliases needed.
\newcommand{\discussion}[1]{\discussionquestion{#1}}

% Simple divider for slide sections
\newcommand{\sectiondivider}{%
    \begin{center}
    \rule{0.8\textwidth}{0.4pt}
    \end{center}
}

% Argument environment (for standard form arguments)
\newenvironment{argument}{%
  \begin{argumentbox}[Argument in Standard Form]
  \begin{enumerate}
}{%
  \end{enumerate}
  \end{argumentbox}
}

% Quote slide environment (combines icon and quotebox)
\newenvironment{quoteslide}[2]{%
  \begin{quotebox}[#1]
  \small
  \textit{``#2''}
}{%
  \end{quotebox}
}

% ==============================================================================
% End of Common Preamble
% ==============================================================================

%
% Individual lecture files should:
% 1. Declare \documentclass (beamer or beamerarticle)
% 2. % ==============================================================================
% Common Preamble for IT Ethics Lectures
% ==============================================================================
% This file contains shared configuration for all lecture files.
% Include this after \documentclass using: \input{lecture_preamble.tex}
%
% Individual lecture files should:
% 1. Declare \documentclass (beamer or beamerarticle)
% 2. \input{lecture_preamble.tex}
% 3. Define lecture-specific colors/customizations
% 4. Set title information
% ==============================================================================

% ------------------------------------------------------------------------------
% Theme and Basic Setup
% ------------------------------------------------------------------------------
\usetheme[progressbar=foot, block=fill, sectionpage=progressbar]{metropolis}

% ------------------------------------------------------------------------------
% Core Packages
% ------------------------------------------------------------------------------
\usepackage{tikz}
\usetikzlibrary{shapes,arrows,positioning,calc,decorations.pathreplacing,fit,backgrounds,chains}
\usepackage{booktabs}
\usepackage{array}
\usepackage{graphicx}
\usepackage{xcolor}
\usepackage[sfdefault]{plex-sans}   % IBM Plex Sans — compact, modern, tech-themed
\usepackage[T1]{fontenc}
\usepackage{fontawesome5}          % real vector icons

% Conditional text boxes (more flexible than basic beamer blocks)
\usepackage{tcolorbox}
\tcbuselibrary{skins,breakable}

% ------------------------------------------------------------------------------
% Bibliography (loaded in all modes for citations)
% ------------------------------------------------------------------------------
\usepackage[style=authoryear,backend=biber,natbib=true,maxcitenames=2]{biblatex}
\addbibresource{refs.bib}

% ------------------------------------------------------------------------------
% Article Mode Adjustments
% ------------------------------------------------------------------------------
% When compiling as article (using beamerarticle class), apply these settings
\mode<article>{
    % Use standard article sections instead of beamer frames
    \usepackage{fullpage}
    \usepackage{hyperref}
    \hypersetup{
        colorlinks=true,
        linkcolor=blue,
        citecolor=blue,
        urlcolor=blue
    }

    % Redefine frame environment for article mode
    % Frames become subsections in article mode
    \renewenvironment{frame}[1][]{
        \par\medskip
        \noindent\textbf{#1}\par\medskip
    }{}
}

% Presentation mode specific settings
\mode<presentation>{
    \setlength{\parskip}{0pt}

    % Show a TOC slide at the start of each section, current section highlighted
    \AtBeginSection[]{
        \begin{frame}{Lecture Outline}
            \tableofcontents[currentsection, hideallsubsections]
        \end{frame}
    }

    % FiraSans has ~15% taller metrics than Computer Modern; pull line spacing back
    % to avoid the massive overflows Metropolis + FiraSans produce by default.
    \renewcommand{\baselinestretch}{0.88}

    % Tighter list spacing
    \setbeamertemplate{itemize/enumerate body begin}{\setlength{\itemsep}{2pt}\setlength{\parsep}{0pt}}
    \setbeamertemplate{itemize/enumerate subbody begin}{\setlength{\itemsep}{1pt}\setlength{\parsep}{0pt}}
}

% ------------------------------------------------------------------------------
% Standard Colors for Boxes
% ------------------------------------------------------------------------------
% Consistent color scheme across all lectures for semantic elements
\definecolor{boxblue}{HTML}{1A4A7A}       % Arguments, concepts
\definecolor{boxgreen}{HTML}{2E7D32}      % Pro arguments, supporting points
\definecolor{boxred}{HTML}{C0392B}        % Objections, warnings
\definecolor{boxgold}{HTML}{B7770D}       % Case studies, examples
\definecolor{boxgray}{HTML}{555F6D}       % Quotes, citations
\definecolor{boxpurple}{HTML}{6A3093}     % Secondary accent
% Metropolis structural colors
\definecolor{mPrimary}{HTML}{1E3A5F}      % Deep navy — frametitle / title bg
\definecolor{mAccent}{HTML}{D4880A}       % Warm amber — progress bar / separators

% ------------------------------------------------------------------------------
% Beamer Theme Color Setup
% ------------------------------------------------------------------------------
% Applied once here; individual lectures should NOT override these.
\mode<presentation>{
    % Metropolis palette primary drives frametitle + title frame background
    \setbeamercolor{palette primary}{bg=mPrimary, fg=white}
    \setbeamercolor{frametitle}{bg=mPrimary, fg=white}
    % Title page element colors — must be explicit so they stay white
    % when a \usebackgroundtemplate image replaces the default navy fill
    \setbeamercolor{title}{fg=white}
    \setbeamercolor{subtitle}{fg=mAccent}
    \setbeamercolor{author}{fg=white}
    \setbeamercolor{institute}{fg=white!85!black}
    \setbeamercolor{date}{fg=white!75!black}
    % Progress bar and title separator use accent color
    \setbeamercolor{progress bar}{fg=mAccent, bg=mAccent!25}
    \setbeamercolor{title separator}{fg=mAccent}
    \setbeamercolor{alerted text}{fg=boxred}
    % Standard beamer blocks
    \setbeamercolor{block title}{fg=white, bg=boxblue}
    \setbeamercolor{block body}{bg=boxblue!8}
    \setbeamercolor{block title alerted}{fg=white, bg=boxred}
    \setbeamercolor{block body alerted}{bg=boxred!8}
    \setbeamercolor{block title example}{fg=white, bg=boxgreen}
    \setbeamercolor{block body example}{bg=boxgreen!5}
}

% ------------------------------------------------------------------------------
% Standard Custom Boxes
% ------------------------------------------------------------------------------
% Uniform boxes used across all lectures - do not redefine in individual files

% Argument/Reasoning Box (logical arguments in standard form)
\newtcolorbox{argumentbox}[1][Argument]{
    colback=boxblue!5,
    colframe=boxblue,
    fonttitle=\bfseries,
    title={#1},
    boxrule=1pt,
    arc=2pt
}

% Quotation Box (for philosophical quotes, sources)
\newtcolorbox{quotebox}[1][Quote]{
    colback=gray!10,
    colframe=boxgray,
    fonttitle=\itshape,
    title={#1},
    boxrule=0.5pt,
    arc=2pt,
    left=10pt,
    right=10pt,
    leftrule=3pt
}

% Case Study/Example Box
\newtcolorbox{casestudybox}[1][Case Study]{
    colback=boxgold!10,
    colframe=boxgold,
    fonttitle=\bfseries,
    title={#1},
    boxrule=1pt,
    arc=2pt
}

% Objection/Counterargument Box
\newtcolorbox{objectionbox}[1][Objection]{
    colback=boxred!5,
    colframe=boxred,
    fonttitle=\bfseries,
    title={#1},
    boxrule=1pt,
    arc=2pt
}

% Pro/Supporting Argument Box
\newtcolorbox{probox}[1][Supporting Argument]{
    colback=boxgreen!5,
    colframe=boxgreen,
    fonttitle=\bfseries,
    title={#1},
    boxrule=1pt,
    arc=2pt
}

% Concept/Definition Box
\newtcolorbox{conceptbox}[1][Key Concept]{
    colback=boxblue!5,
    colframe=boxblue,
    fonttitle=\bfseries,
    title={#1},
    boxrule=1pt,
    arc=2pt
}

% Discussion Box (for group activities, questions)
\newtcolorbox{discussionbox}[1][Discussion]{
    colback=boxgold!5,
    colframe=boxgold,
    fonttitle=\bfseries,
    title={#1},
    boxrule=1pt,
    arc=2pt
}

% Figure Box (for diagrams with explanations)
\newtcolorbox{figurebox}[1][Figure]{
    colback=boxblue!3,
    colframe=boxblue!70,
    fonttitle=\bfseries,
    title={#1},
    boxrule=0.75pt,
    arc=2pt
}

% ------------------------------------------------------------------------------
% Utility Commands and Icons
% ------------------------------------------------------------------------------

% fontawesome5 icon shorthands (\faXxx commands come from the fontawesome5 package)
\newcommand{\iconQuote}{\faQuoteLeft}
\newcommand{\iconBalance}{\faBalanceScale}
\newcommand{\iconDiscussion}{\faComments}
\newcommand{\iconIdea}{\faLightbulb}
\newcommand{\iconPro}{\faThumbsUp}
\newcommand{\iconCon}{\faThumbsDown}

% Discussion question command (for interactive moments)
\newcommand{\discussionquestion}[1]{%
    \begin{alertblock}{Discussion Question}
    #1
    \end{alertblock}
}

% fontawesome5 defines \faQuoteLeft, \faBalanceScale, etc. natively --- no aliases needed.
\newcommand{\discussion}[1]{\discussionquestion{#1}}

% Simple divider for slide sections
\newcommand{\sectiondivider}{%
    \begin{center}
    \rule{0.8\textwidth}{0.4pt}
    \end{center}
}

% Argument environment (for standard form arguments)
\newenvironment{argument}{%
  \begin{argumentbox}[Argument in Standard Form]
  \begin{enumerate}
}{%
  \end{enumerate}
  \end{argumentbox}
}

% Quote slide environment (combines icon and quotebox)
\newenvironment{quoteslide}[2]{%
  \begin{quotebox}[#1]
  \small
  \textit{``#2''}
}{%
  \end{quotebox}
}

% ==============================================================================
% End of Common Preamble
% ==============================================================================

% 3. Define lecture-specific colors/customizations
% 4. Set title information
% ==============================================================================

% ------------------------------------------------------------------------------
% Theme and Basic Setup
% ------------------------------------------------------------------------------
\usetheme[progressbar=foot, block=fill, sectionpage=progressbar]{metropolis}

% ------------------------------------------------------------------------------
% Core Packages
% ------------------------------------------------------------------------------
\usepackage{tikz}
\usetikzlibrary{shapes,arrows,positioning,calc,decorations.pathreplacing,fit,backgrounds,chains}
\usepackage{booktabs}
\usepackage{array}
\usepackage{graphicx}
\usepackage{xcolor}
\usepackage[sfdefault]{plex-sans}   % IBM Plex Sans — compact, modern, tech-themed
\usepackage[T1]{fontenc}
\usepackage{fontawesome5}          % real vector icons

% Conditional text boxes (more flexible than basic beamer blocks)
\usepackage{tcolorbox}
\tcbuselibrary{skins,breakable}

% ------------------------------------------------------------------------------
% Bibliography (loaded in all modes for citations)
% ------------------------------------------------------------------------------
\usepackage[style=authoryear,backend=biber,natbib=true,maxcitenames=2]{biblatex}
\addbibresource{refs.bib}

% ------------------------------------------------------------------------------
% Article Mode Adjustments
% ------------------------------------------------------------------------------
% When compiling as article (using beamerarticle class), apply these settings
\mode<article>{
    % Use standard article sections instead of beamer frames
    \usepackage{fullpage}
    \usepackage{hyperref}
    \hypersetup{
        colorlinks=true,
        linkcolor=blue,
        citecolor=blue,
        urlcolor=blue
    }

    % Redefine frame environment for article mode
    % Frames become subsections in article mode
    \renewenvironment{frame}[1][]{
        \par\medskip
        \noindent\textbf{#1}\par\medskip
    }{}
}

% Presentation mode specific settings
\mode<presentation>{
    \setlength{\parskip}{0pt}

    % Show a TOC slide at the start of each section, current section highlighted
    \AtBeginSection[]{
        \begin{frame}{Lecture Outline}
            \tableofcontents[currentsection, hideallsubsections]
        \end{frame}
    }

    % FiraSans has ~15% taller metrics than Computer Modern; pull line spacing back
    % to avoid the massive overflows Metropolis + FiraSans produce by default.
    \renewcommand{\baselinestretch}{0.88}

    % Tighter list spacing
    \setbeamertemplate{itemize/enumerate body begin}{\setlength{\itemsep}{2pt}\setlength{\parsep}{0pt}}
    \setbeamertemplate{itemize/enumerate subbody begin}{\setlength{\itemsep}{1pt}\setlength{\parsep}{0pt}}
}

% ------------------------------------------------------------------------------
% Standard Colors for Boxes
% ------------------------------------------------------------------------------
% Consistent color scheme across all lectures for semantic elements
\definecolor{boxblue}{HTML}{1A4A7A}       % Arguments, concepts
\definecolor{boxgreen}{HTML}{2E7D32}      % Pro arguments, supporting points
\definecolor{boxred}{HTML}{C0392B}        % Objections, warnings
\definecolor{boxgold}{HTML}{B7770D}       % Case studies, examples
\definecolor{boxgray}{HTML}{555F6D}       % Quotes, citations
\definecolor{boxpurple}{HTML}{6A3093}     % Secondary accent
% Metropolis structural colors
\definecolor{mPrimary}{HTML}{1E3A5F}      % Deep navy — frametitle / title bg
\definecolor{mAccent}{HTML}{D4880A}       % Warm amber — progress bar / separators

% ------------------------------------------------------------------------------
% Beamer Theme Color Setup
% ------------------------------------------------------------------------------
% Applied once here; individual lectures should NOT override these.
\mode<presentation>{
    % Metropolis palette primary drives frametitle + title frame background
    \setbeamercolor{palette primary}{bg=mPrimary, fg=white}
    \setbeamercolor{frametitle}{bg=mPrimary, fg=white}
    % Title page element colors — must be explicit so they stay white
    % when a \usebackgroundtemplate image replaces the default navy fill
    \setbeamercolor{title}{fg=white}
    \setbeamercolor{subtitle}{fg=mAccent}
    \setbeamercolor{author}{fg=white}
    \setbeamercolor{institute}{fg=white!85!black}
    \setbeamercolor{date}{fg=white!75!black}
    % Progress bar and title separator use accent color
    \setbeamercolor{progress bar}{fg=mAccent, bg=mAccent!25}
    \setbeamercolor{title separator}{fg=mAccent}
    \setbeamercolor{alerted text}{fg=boxred}
    % Standard beamer blocks
    \setbeamercolor{block title}{fg=white, bg=boxblue}
    \setbeamercolor{block body}{bg=boxblue!8}
    \setbeamercolor{block title alerted}{fg=white, bg=boxred}
    \setbeamercolor{block body alerted}{bg=boxred!8}
    \setbeamercolor{block title example}{fg=white, bg=boxgreen}
    \setbeamercolor{block body example}{bg=boxgreen!5}
}

% ------------------------------------------------------------------------------
% Standard Custom Boxes
% ------------------------------------------------------------------------------
% Uniform boxes used across all lectures - do not redefine in individual files

% Argument/Reasoning Box (logical arguments in standard form)
\newtcolorbox{argumentbox}[1][Argument]{
    colback=boxblue!5,
    colframe=boxblue,
    fonttitle=\bfseries,
    title={#1},
    boxrule=1pt,
    arc=2pt
}

% Quotation Box (for philosophical quotes, sources)
\newtcolorbox{quotebox}[1][Quote]{
    colback=gray!10,
    colframe=boxgray,
    fonttitle=\itshape,
    title={#1},
    boxrule=0.5pt,
    arc=2pt,
    left=10pt,
    right=10pt,
    leftrule=3pt
}

% Case Study/Example Box
\newtcolorbox{casestudybox}[1][Case Study]{
    colback=boxgold!10,
    colframe=boxgold,
    fonttitle=\bfseries,
    title={#1},
    boxrule=1pt,
    arc=2pt
}

% Objection/Counterargument Box
\newtcolorbox{objectionbox}[1][Objection]{
    colback=boxred!5,
    colframe=boxred,
    fonttitle=\bfseries,
    title={#1},
    boxrule=1pt,
    arc=2pt
}

% Pro/Supporting Argument Box
\newtcolorbox{probox}[1][Supporting Argument]{
    colback=boxgreen!5,
    colframe=boxgreen,
    fonttitle=\bfseries,
    title={#1},
    boxrule=1pt,
    arc=2pt
}

% Concept/Definition Box
\newtcolorbox{conceptbox}[1][Key Concept]{
    colback=boxblue!5,
    colframe=boxblue,
    fonttitle=\bfseries,
    title={#1},
    boxrule=1pt,
    arc=2pt
}

% Discussion Box (for group activities, questions)
\newtcolorbox{discussionbox}[1][Discussion]{
    colback=boxgold!5,
    colframe=boxgold,
    fonttitle=\bfseries,
    title={#1},
    boxrule=1pt,
    arc=2pt
}

% Figure Box (for diagrams with explanations)
\newtcolorbox{figurebox}[1][Figure]{
    colback=boxblue!3,
    colframe=boxblue!70,
    fonttitle=\bfseries,
    title={#1},
    boxrule=0.75pt,
    arc=2pt
}

% ------------------------------------------------------------------------------
% Utility Commands and Icons
% ------------------------------------------------------------------------------

% fontawesome5 icon shorthands (\faXxx commands come from the fontawesome5 package)
\newcommand{\iconQuote}{\faQuoteLeft}
\newcommand{\iconBalance}{\faBalanceScale}
\newcommand{\iconDiscussion}{\faComments}
\newcommand{\iconIdea}{\faLightbulb}
\newcommand{\iconPro}{\faThumbsUp}
\newcommand{\iconCon}{\faThumbsDown}

% Discussion question command (for interactive moments)
\newcommand{\discussionquestion}[1]{%
    \begin{alertblock}{Discussion Question}
    #1
    \end{alertblock}
}

% fontawesome5 defines \faQuoteLeft, \faBalanceScale, etc. natively --- no aliases needed.
\newcommand{\discussion}[1]{\discussionquestion{#1}}

% Simple divider for slide sections
\newcommand{\sectiondivider}{%
    \begin{center}
    \rule{0.8\textwidth}{0.4pt}
    \end{center}
}

% Argument environment (for standard form arguments)
\newenvironment{argument}{%
  \begin{argumentbox}[Argument in Standard Form]
  \begin{enumerate}
}{%
  \end{enumerate}
  \end{argumentbox}
}

% Quote slide environment (combines icon and quotebox)
\newenvironment{quoteslide}[2]{%
  \begin{quotebox}[#1]
  \small
  \textit{``#2''}
}{%
  \end{quotebox}
}

% ==============================================================================
% End of Common Preamble
% ==============================================================================

% 3. Define lecture-specific colors/customizations
% 4. Set title information
% ==============================================================================

% ------------------------------------------------------------------------------
% Theme and Basic Setup
% ------------------------------------------------------------------------------
\usetheme[progressbar=foot, block=fill, sectionpage=progressbar]{metropolis}

% ------------------------------------------------------------------------------
% Core Packages
% ------------------------------------------------------------------------------
\usepackage{tikz}
\usetikzlibrary{shapes,arrows,positioning,calc,decorations.pathreplacing,fit,backgrounds,chains}
\usepackage{booktabs}
\usepackage{array}
\usepackage{graphicx}
\usepackage{xcolor}
\usepackage[sfdefault]{plex-sans}   % IBM Plex Sans — compact, modern, tech-themed
\usepackage[T1]{fontenc}
\usepackage{fontawesome5}          % real vector icons

% Conditional text boxes (more flexible than basic beamer blocks)
\usepackage{tcolorbox}
\tcbuselibrary{skins,breakable}

% ------------------------------------------------------------------------------
% Bibliography (loaded in all modes for citations)
% ------------------------------------------------------------------------------
\usepackage[style=authoryear,backend=biber,natbib=true,maxcitenames=2]{biblatex}
\addbibresource{refs.bib}

% ------------------------------------------------------------------------------
% Article Mode Adjustments
% ------------------------------------------------------------------------------
% When compiling as article (using beamerarticle class), apply these settings
\mode<article>{
    % Use standard article sections instead of beamer frames
    \usepackage{fullpage}
    \usepackage{hyperref}
    \hypersetup{
        colorlinks=true,
        linkcolor=blue,
        citecolor=blue,
        urlcolor=blue
    }

    % Redefine frame environment for article mode
    % Frames become subsections in article mode
    \renewenvironment{frame}[1][]{
        \par\medskip
        \noindent\textbf{#1}\par\medskip
    }{}
}

% Presentation mode specific settings
\mode<presentation>{
    \setlength{\parskip}{0pt}

    % Show a TOC slide at the start of each section, current section highlighted
    \AtBeginSection[]{
        \begin{frame}{Lecture Outline}
            \tableofcontents[currentsection, hideallsubsections]
        \end{frame}
    }

    % FiraSans has ~15% taller metrics than Computer Modern; pull line spacing back
    % to avoid the massive overflows Metropolis + FiraSans produce by default.
    \renewcommand{\baselinestretch}{0.88}

    % Tighter list spacing
    \setbeamertemplate{itemize/enumerate body begin}{\setlength{\itemsep}{2pt}\setlength{\parsep}{0pt}}
    \setbeamertemplate{itemize/enumerate subbody begin}{\setlength{\itemsep}{1pt}\setlength{\parsep}{0pt}}
}

% ------------------------------------------------------------------------------
% Standard Colors for Boxes
% ------------------------------------------------------------------------------
% Consistent color scheme across all lectures for semantic elements
\definecolor{boxblue}{HTML}{1A4A7A}       % Arguments, concepts
\definecolor{boxgreen}{HTML}{2E7D32}      % Pro arguments, supporting points
\definecolor{boxred}{HTML}{C0392B}        % Objections, warnings
\definecolor{boxgold}{HTML}{B7770D}       % Case studies, examples
\definecolor{boxgray}{HTML}{555F6D}       % Quotes, citations
\definecolor{boxpurple}{HTML}{6A3093}     % Secondary accent
% Metropolis structural colors
\definecolor{mPrimary}{HTML}{1E3A5F}      % Deep navy — frametitle / title bg
\definecolor{mAccent}{HTML}{D4880A}       % Warm amber — progress bar / separators

% ------------------------------------------------------------------------------
% Beamer Theme Color Setup
% ------------------------------------------------------------------------------
% Applied once here; individual lectures should NOT override these.
\mode<presentation>{
    % Metropolis palette primary drives frametitle + title frame background
    \setbeamercolor{palette primary}{bg=mPrimary, fg=white}
    \setbeamercolor{frametitle}{bg=mPrimary, fg=white}
    % Title page element colors — must be explicit so they stay white
    % when a \usebackgroundtemplate image replaces the default navy fill
    \setbeamercolor{title}{fg=white}
    \setbeamercolor{subtitle}{fg=mAccent}
    \setbeamercolor{author}{fg=white}
    \setbeamercolor{institute}{fg=white!85!black}
    \setbeamercolor{date}{fg=white!75!black}
    % Progress bar and title separator use accent color
    \setbeamercolor{progress bar}{fg=mAccent, bg=mAccent!25}
    \setbeamercolor{title separator}{fg=mAccent}
    \setbeamercolor{alerted text}{fg=boxred}
    % Standard beamer blocks
    \setbeamercolor{block title}{fg=white, bg=boxblue}
    \setbeamercolor{block body}{bg=boxblue!8}
    \setbeamercolor{block title alerted}{fg=white, bg=boxred}
    \setbeamercolor{block body alerted}{bg=boxred!8}
    \setbeamercolor{block title example}{fg=white, bg=boxgreen}
    \setbeamercolor{block body example}{bg=boxgreen!5}
}

% ------------------------------------------------------------------------------
% Standard Custom Boxes
% ------------------------------------------------------------------------------
% Uniform boxes used across all lectures - do not redefine in individual files

% Argument/Reasoning Box (logical arguments in standard form)
\newtcolorbox{argumentbox}[1][Argument]{
    colback=boxblue!5,
    colframe=boxblue,
    fonttitle=\bfseries,
    title={#1},
    boxrule=1pt,
    arc=2pt
}

% Quotation Box (for philosophical quotes, sources)
\newtcolorbox{quotebox}[1][Quote]{
    colback=gray!10,
    colframe=boxgray,
    fonttitle=\itshape,
    title={#1},
    boxrule=0.5pt,
    arc=2pt,
    left=10pt,
    right=10pt,
    leftrule=3pt
}

% Case Study/Example Box
\newtcolorbox{casestudybox}[1][Case Study]{
    colback=boxgold!10,
    colframe=boxgold,
    fonttitle=\bfseries,
    title={#1},
    boxrule=1pt,
    arc=2pt
}

% Objection/Counterargument Box
\newtcolorbox{objectionbox}[1][Objection]{
    colback=boxred!5,
    colframe=boxred,
    fonttitle=\bfseries,
    title={#1},
    boxrule=1pt,
    arc=2pt
}

% Pro/Supporting Argument Box
\newtcolorbox{probox}[1][Supporting Argument]{
    colback=boxgreen!5,
    colframe=boxgreen,
    fonttitle=\bfseries,
    title={#1},
    boxrule=1pt,
    arc=2pt
}

% Concept/Definition Box
\newtcolorbox{conceptbox}[1][Key Concept]{
    colback=boxblue!5,
    colframe=boxblue,
    fonttitle=\bfseries,
    title={#1},
    boxrule=1pt,
    arc=2pt
}

% Discussion Box (for group activities, questions)
\newtcolorbox{discussionbox}[1][Discussion]{
    colback=boxgold!5,
    colframe=boxgold,
    fonttitle=\bfseries,
    title={#1},
    boxrule=1pt,
    arc=2pt
}

% Figure Box (for diagrams with explanations)
\newtcolorbox{figurebox}[1][Figure]{
    colback=boxblue!3,
    colframe=boxblue!70,
    fonttitle=\bfseries,
    title={#1},
    boxrule=0.75pt,
    arc=2pt
}

% ------------------------------------------------------------------------------
% Utility Commands and Icons
% ------------------------------------------------------------------------------

% fontawesome5 icon shorthands (\faXxx commands come from the fontawesome5 package)
\newcommand{\iconQuote}{\faQuoteLeft}
\newcommand{\iconBalance}{\faBalanceScale}
\newcommand{\iconDiscussion}{\faComments}
\newcommand{\iconIdea}{\faLightbulb}
\newcommand{\iconPro}{\faThumbsUp}
\newcommand{\iconCon}{\faThumbsDown}

% Discussion question command (for interactive moments)
\newcommand{\discussionquestion}[1]{%
    \begin{alertblock}{Discussion Question}
    #1
    \end{alertblock}
}

% fontawesome5 defines \faQuoteLeft, \faBalanceScale, etc. natively --- no aliases needed.
\newcommand{\discussion}[1]{\discussionquestion{#1}}

% Simple divider for slide sections
\newcommand{\sectiondivider}{%
    \begin{center}
    \rule{0.8\textwidth}{0.4pt}
    \end{center}
}

% Argument environment (for standard form arguments)
\newenvironment{argument}{%
  \begin{argumentbox}[Argument in Standard Form]
  \begin{enumerate}
}{%
  \end{enumerate}
  \end{argumentbox}
}

% Quote slide environment (combines icon and quotebox)
\newenvironment{quoteslide}[2]{%
  \begin{quotebox}[#1]
  \small
  \textit{``#2''}
}{%
  \end{quotebox}
}

% ==============================================================================
% End of Common Preamble
% ==============================================================================

% 3. Define lecture-specific colors/customizations
% 4. Set title information
% ==============================================================================

% ------------------------------------------------------------------------------
% Theme and Basic Setup
% ------------------------------------------------------------------------------
\usetheme[progressbar=foot, block=fill, sectionpage=progressbar]{metropolis}

% ------------------------------------------------------------------------------
% Core Packages
% ------------------------------------------------------------------------------
\usepackage{tikz}
\usetikzlibrary{shapes,arrows,positioning,calc,decorations.pathreplacing,fit,backgrounds,chains}
\usepackage{booktabs}
\usepackage{array}
\usepackage{graphicx}
\usepackage{xcolor}
\usepackage[sfdefault]{plex-sans}   % IBM Plex Sans — compact, modern, tech-themed
\usepackage[T1]{fontenc}
\usepackage{fontawesome5}          % real vector icons

% Conditional text boxes (more flexible than basic beamer blocks)
\usepackage{tcolorbox}
\tcbuselibrary{skins,breakable}

% ------------------------------------------------------------------------------
% Bibliography (loaded in all modes for citations)
% ------------------------------------------------------------------------------
\usepackage[style=authoryear,backend=biber,natbib=true,maxcitenames=2]{biblatex}
\addbibresource{refs.bib}

% ------------------------------------------------------------------------------
% Article Mode Adjustments
% ------------------------------------------------------------------------------
% When compiling as article (using beamerarticle class), apply these settings
\mode<article>{
    % Use standard article sections instead of beamer frames
    \usepackage{fullpage}
    \usepackage{hyperref}
    \hypersetup{
        colorlinks=true,
        linkcolor=blue,
        citecolor=blue,
        urlcolor=blue
    }

    % Redefine frame environment for article mode
    % Frames become subsections in article mode
    \renewenvironment{frame}[1][]{
        \par\medskip
        \noindent\textbf{#1}\par\medskip
    }{}
}

% Presentation mode specific settings
\mode<presentation>{
    \setlength{\parskip}{0pt}

    % Show a TOC slide at the start of each section, current section highlighted
    \AtBeginSection[]{
        \begin{frame}{Lecture Outline}
            \tableofcontents[currentsection, hideallsubsections]
        \end{frame}
    }

    % FiraSans has ~15% taller metrics than Computer Modern; pull line spacing back
    % to avoid the massive overflows Metropolis + FiraSans produce by default.
    \renewcommand{\baselinestretch}{0.88}

    % Tighter list spacing
    \setbeamertemplate{itemize/enumerate body begin}{\setlength{\itemsep}{2pt}\setlength{\parsep}{0pt}}
    \setbeamertemplate{itemize/enumerate subbody begin}{\setlength{\itemsep}{1pt}\setlength{\parsep}{0pt}}
}

% ------------------------------------------------------------------------------
% Standard Colors for Boxes
% ------------------------------------------------------------------------------
% Consistent color scheme across all lectures for semantic elements
\definecolor{boxblue}{HTML}{1A4A7A}       % Arguments, concepts
\definecolor{boxgreen}{HTML}{2E7D32}      % Pro arguments, supporting points
\definecolor{boxred}{HTML}{C0392B}        % Objections, warnings
\definecolor{boxgold}{HTML}{B7770D}       % Case studies, examples
\definecolor{boxgray}{HTML}{555F6D}       % Quotes, citations
\definecolor{boxpurple}{HTML}{6A3093}     % Secondary accent
% Metropolis structural colors
\definecolor{mPrimary}{HTML}{1E3A5F}      % Deep navy — frametitle / title bg
\definecolor{mAccent}{HTML}{D4880A}       % Warm amber — progress bar / separators

% ------------------------------------------------------------------------------
% Beamer Theme Color Setup
% ------------------------------------------------------------------------------
% Applied once here; individual lectures should NOT override these.
\mode<presentation>{
    % Metropolis palette primary drives frametitle + title frame background
    \setbeamercolor{palette primary}{bg=mPrimary, fg=white}
    \setbeamercolor{frametitle}{bg=mPrimary, fg=white}
    % Title page element colors — must be explicit so they stay white
    % when a \usebackgroundtemplate image replaces the default navy fill
    \setbeamercolor{title}{fg=white}
    \setbeamercolor{subtitle}{fg=mAccent}
    \setbeamercolor{author}{fg=white}
    \setbeamercolor{institute}{fg=white!85!black}
    \setbeamercolor{date}{fg=white!75!black}
    % Progress bar and title separator use accent color
    \setbeamercolor{progress bar}{fg=mAccent, bg=mAccent!25}
    \setbeamercolor{title separator}{fg=mAccent}
    \setbeamercolor{alerted text}{fg=boxred}
    % Standard beamer blocks
    \setbeamercolor{block title}{fg=white, bg=boxblue}
    \setbeamercolor{block body}{bg=boxblue!8}
    \setbeamercolor{block title alerted}{fg=white, bg=boxred}
    \setbeamercolor{block body alerted}{bg=boxred!8}
    \setbeamercolor{block title example}{fg=white, bg=boxgreen}
    \setbeamercolor{block body example}{bg=boxgreen!5}
}

% ------------------------------------------------------------------------------
% Standard Custom Boxes
% ------------------------------------------------------------------------------
% Uniform boxes used across all lectures - do not redefine in individual files

% Argument/Reasoning Box (logical arguments in standard form)
\newtcolorbox{argumentbox}[1][Argument]{
    colback=boxblue!5,
    colframe=boxblue,
    fonttitle=\bfseries,
    title={#1},
    boxrule=1pt,
    arc=2pt
}

% Quotation Box (for philosophical quotes, sources)
\newtcolorbox{quotebox}[1][Quote]{
    colback=gray!10,
    colframe=boxgray,
    fonttitle=\itshape,
    title={#1},
    boxrule=0.5pt,
    arc=2pt,
    left=10pt,
    right=10pt,
    leftrule=3pt
}

% Case Study/Example Box
\newtcolorbox{casestudybox}[1][Case Study]{
    colback=boxgold!10,
    colframe=boxgold,
    fonttitle=\bfseries,
    title={#1},
    boxrule=1pt,
    arc=2pt
}

% Objection/Counterargument Box
\newtcolorbox{objectionbox}[1][Objection]{
    colback=boxred!5,
    colframe=boxred,
    fonttitle=\bfseries,
    title={#1},
    boxrule=1pt,
    arc=2pt
}

% Pro/Supporting Argument Box
\newtcolorbox{probox}[1][Supporting Argument]{
    colback=boxgreen!5,
    colframe=boxgreen,
    fonttitle=\bfseries,
    title={#1},
    boxrule=1pt,
    arc=2pt
}

% Concept/Definition Box
\newtcolorbox{conceptbox}[1][Key Concept]{
    colback=boxblue!5,
    colframe=boxblue,
    fonttitle=\bfseries,
    title={#1},
    boxrule=1pt,
    arc=2pt
}

% Discussion Box (for group activities, questions)
\newtcolorbox{discussionbox}[1][Discussion]{
    colback=boxgold!5,
    colframe=boxgold,
    fonttitle=\bfseries,
    title={#1},
    boxrule=1pt,
    arc=2pt
}

% Figure Box (for diagrams with explanations)
\newtcolorbox{figurebox}[1][Figure]{
    colback=boxblue!3,
    colframe=boxblue!70,
    fonttitle=\bfseries,
    title={#1},
    boxrule=0.75pt,
    arc=2pt
}

% ------------------------------------------------------------------------------
% Utility Commands and Icons
% ------------------------------------------------------------------------------

% fontawesome5 icon shorthands (\faXxx commands come from the fontawesome5 package)
\newcommand{\iconQuote}{\faQuoteLeft}
\newcommand{\iconBalance}{\faBalanceScale}
\newcommand{\iconDiscussion}{\faComments}
\newcommand{\iconIdea}{\faLightbulb}
\newcommand{\iconPro}{\faThumbsUp}
\newcommand{\iconCon}{\faThumbsDown}

% Discussion question command (for interactive moments)
\newcommand{\discussionquestion}[1]{%
    \begin{alertblock}{Discussion Question}
    #1
    \end{alertblock}
}

% fontawesome5 defines \faQuoteLeft, \faBalanceScale, etc. natively --- no aliases needed.
\newcommand{\discussion}[1]{\discussionquestion{#1}}

% Simple divider for slide sections
\newcommand{\sectiondivider}{%
    \begin{center}
    \rule{0.8\textwidth}{0.4pt}
    \end{center}
}

% Argument environment (for standard form arguments)
\newenvironment{argument}{%
  \begin{argumentbox}[Argument in Standard Form]
  \begin{enumerate}
}{%
  \end{enumerate}
  \end{argumentbox}
}

% Quote slide environment (combines icon and quotebox)
\newenvironment{quoteslide}[2]{%
  \begin{quotebox}[#1]
  \small
  \textit{``#2''}
}{%
  \end{quotebox}
}

% ==============================================================================
% End of Common Preamble
% ==============================================================================
